% =====================================================================
%  Chapter 6: Hadron spectroscopy
%  Gattringer & Lang, "Quantum Chromodynamics on the Lattice" (LNP 788)
%  Reconstructed from extract/ch06.txt; equations auto-numbered.
% =====================================================================

\chapter{Hadron spectroscopy}

The simplest quantities involving fermions that one can compute on the lattice
are the masses of hadrons. Since there are many combinations of quantum
numbers such as spin, parity, flavor content etc.\ there is a wealth of
possible hadrons. Reproducing all their masses correctly is already a powerful
test for the correctness of QCD. Over the last 20 years lattice QCD
calculations of the mass spectrum have improved continually and, for ground
state baryons, have reached impressive agreement with experimental data.

In this chapter we discuss the conceptual background and the actual
implementation of a numerical spectroscopy calculation. The hadron correlation
functions are not only necessary for computing the spectrum but enter also in
the determination of hadronic matrix elements addressed in later chapters. We
first discuss how to construct operators with the correct quantum numbers and
their correlation functions and address the quenched approximation.
Subsequently we present techniques for quark sources and show how quark
propagators can be computed. We continue with the analysis of the resulting
hadron propagators and discuss how to obtain the hadron masses.

\section{Hadron interpolators and correlators}

In a typical QCD calculation one generates gauge configurations according to
the requested distribution and then computes various observables. These
include simple expectation values of plaquettes or of more extended operators
such as Wilson loops. Also quantities like the chiral condensate
$\langle\bar\psi\psi\rangle$ have been studied and will be discussed in later
chapters.

Here we concentrate on determining hadron masses, i.e., we perform a hadron
spectroscopy calculation. For this purpose one first computes quark
propagators for each gauge configuration. These are then suitably combined to
construct the hadron propagators. Finally, averaging over all gauge
configurations provides our estimate of the hadron propagators for subsequent
analysis.

The first step of such a spectroscopy calculation is the identification of
hadron interpolators $\bar{O}$, $O$ such that the corresponding Hilbert space
operators $\hat{\bar{O}}$, $\hat{O}^\dagger$ annihilate and create the particle
states we want to analyze. A hadron interpolator is a functional of the
lattice fields with the quantum numbers of the state one is interested
in.\footnote{According to common practice, we will often use ``operator'' and
``interpolator'' synonymously in the following.} Once the interpolators are
identified we consider the Euclidean correlator (1.12) of hadron interpolators
$\bar{O}(n_t)$, $O(0)$ located at time slices $n_4 = n_t$ and $n_4 = 0$.

For hadron spectroscopy one studies interpolators constructed out of quarks
and gluons. These are by construction gauge-invariant color singlets. Such
operators include:
\begin{itemize}
  \item Local meson operators, like $O_M(n) \equiv \bar\psi(n)\Gamma\psi(n)$,
        consisting of quarks and antiquarks or baryon operators made out of
        three quarks; these should give information on the mesons and baryons
        as predicted by QCD.
  \item Extended interpolators, also contributing to hadron propagators,
        involving, e.g., for mesons terms like
        $\bar\psi(n)U_\mu(n)\psi(n+\hat\mu)$ and similar for baryons.
  \item Pure gauge field interpolators like the plaquette, longer closed
        loops, or more general gluonic operators. These interpolators couple
        to gluonic objects such as glueballs, but also to mesons.
  \item Other, more exotic color-singlet combinations of quarks and antiquarks
        like $\psi\psi\psi\psi$ and $\psi\psi\bar\psi\bar\psi$.
  \item States with $3n$ quarks; these are just the nuclei. For example, the
        deuterium is a proton--neutron state made out of six valence quarks,
        three of up-type and three of down-type.
\end{itemize}

According to (1.21), physically allowed states can be observed in the spectral
decomposition of the propagators of these interpolators,
\begin{equation}
  \big\langle \bar{O}(n_t)\, O(0) \big\rangle
  = \sum_k \bra{0}\hat{O}\ket{k}\bra{k}\hat{O}^\dagger\ket{0} e^{-n_t a E_k}
  = A\, e^{-n_t a E_H}\left( 1 + \order{e^{-n_t a \Delta E}} \right),
\end{equation}
where $A$ is a constant, $E_H$ is the energy of the lowest state $\ket{H}$
with $\bra{0}\hat{O}\ket{H} \neq 0$, and $\Delta E$ is the energy difference to
the first excited state. From the leading exponential decay we thus can
extract the energy $E_H$ of the hadron.

\subsection{Meson interpolators}

The discrete symmetries which we have discussed in Sect.~5.4 play an important
role in the classification of hadrons and are needed for the construction of
hadron interpolators. To illustrate their use we discuss as an example the
construction of pion interpolators in more detail. Other meson interpolators
can then be constructed along the same lines with the help of Table~6.1, which
lists the quantum numbers of fermion bilinears.

To be specific, we consider QCD with only two or three flavors (up, down,
strange). Instead of using a flavor index $f$ attached to spinors denoted by
$\psi$, $\bar\psi$, we often use a different symbol for each flavor. Thus $u$,
$\bar u$, $d$, $\bar d$ and $s$, $\bar s$ denote the spinors for the three
light flavors.

The pions, which we now consider in our example, are made from $u$ and $d$
quarks only. An up-quark $u$ has isospin $I = 1/2$, $I_z = +1/2$ and charge
$Q = 2/3\,e$, where in our convention the charge of the electron is $-e$. For
the down-quark $d$ we have $I = 1/2$, $I_z = -1/2$ and $Q = -1/3\,e$. The
pseudoscalar combinations can be classified with respect to their isospin and
are grouped into an isotriplet ($I = 1$) containing the particles $\pi^+$,
$\pi^0$, $\pi^-$ with isospin components $I_z = 1, 0, -1$ and an iso-singlet
($I = 0$) containing the $\eta$-meson.\footnote{We remark that the physical
$\eta$ also has an admixture of $s\bar s$. See, e.g., [1].}

The charged pions $\pi^+$ and $\pi^-$ have a mass of 140\,MeV. They have zero
spin ($J = 0$), negative parity ($P = -1$), isospin $I = 1$, $I_z = \pm 1$, and
electric charge $Q = \pm e$. Thus, for obtaining the correct charge and
isospin, $\pi^+$ must be a $\bar d$--$u$ combination and $\pi^-$ is of the type
$u$--$\bar d$. We also expect them to be color singlets. More explicitly we can
use the pseudoscalar interpolators (summation convention for all indices)
\begin{equation}
  \begin{gathered}
    O_{\pi^+}(n) = \bar d(n)\, \gamma_5\, u(n)
      = \bar d(n)^\alpha_c (\gamma_5)_{\alpha\beta} u(n)^\beta_c , \\
    O_{\pi^-}(n) = \bar u(n)\, \gamma_5\, d(n)
      = \bar u(n)^\alpha_c (\gamma_5)_{\alpha\beta} d(n)^\beta_c .
  \end{gathered}
\end{equation}
For their transformation under parity (5.72) we find
\begin{multline}
  O_{\pi^+}(n, n_4) = \bar d(n, n_4)\, \gamma_5\, u(n, n_4)
  \xrightarrow{\;P\;}
  \bar d(-n, n_4)\, \gamma_4\gamma_5\gamma_4\, u(-n, n_4) \\
  = -\bar d(-n, n_4)\, \gamma_5\, u(-n, n_4)
  = -O_{\pi^+}(-n, n_4) ,
\end{multline}
showing that our interpolator $O_{\pi^+}$ indeed has negative parity (the same
holds for $O_{\pi^-}$). The fact that the parity transformation changes the
spatial vector $n$ into $-n$ is irrelevant for interpolators where we project
to zero momentum, i.e., where we sum over all $n$ (see Sect.~6.1.4). Applying
charge conjugation (5.66) gives
\begin{multline}
  O_{\pi^+}(n) = \bar d(n)\, \gamma_5\, u(n)
  \xrightarrow{\;C\;}
  -d(n)^T\, C\gamma_5 C^{-1}\, u(n)^T
  = -d(n)^T\, \gamma_5^T\, u(n)^T \\
  = \bar u(n)\, \gamma_5\, d(n) = O_{\pi^-}(n) ,
\end{multline}
where we used $C\gamma_5 C^{-1} = \gamma_5^T$, which follows from (5.67). The
interchange of the Grassmann variables cancels the overall minus sign. Equation
(6.4) establishes that charge conjugation transforms $O_{\pi^+}$ into
$O_{\pi^-}$ and vice versa. The interpolator for the $I_z = 0$ component of the
iso-triplet, the $\pi^0$, is
\begin{equation}
  O_{\pi^0}(n) = \frac{1}{\sqrt{2}}\left[ \bar u(n)\, \gamma_5\, u(n)
    - \bar d(n)\, \gamma_5\, d(n) \right].
\end{equation}

\begin{table}[htbp]
  \centering
  \caption*{Table 6.1. Quantum numbers of the most commonly used meson
    interpolators according to the general form (6.9). We remark that the
    classification with respect to $C$ is for the flavor neutral interpolators
    only}
  \begin{tabular}{@{}l c c l@{}}
    \toprule
    $J^{PC}$ & State & $\Gamma$ & Particles \\
    \midrule
    $0^{++}$ & Scalar & $1,\ \gamma_4$ & $f_0,\ a_0,\ K_0^*,\ \ldots$ \\
    $0^{-+}$ & Pseudoscalar & $\gamma_5,\ \gamma_4\gamma_5$
      & $\pi^\pm,\ \pi^0,\ \eta,\ K^\pm,\ K^0,\ \ldots$ \\
    $1^{--}$ & Vector & $\gamma_i,\ \gamma_4\gamma_i$
      & $\rho^\pm,\ \rho^0,\ \omega,\ K^*,\ \phi,\ \ldots$ \\
    $1^{++}$ & Axial vector & $\gamma_i\gamma_5$ & $a_1,\ f_1,\ \ldots$ \\
    $1^{+-}$ & Tensor & $\gamma_i\gamma_j$ & $h_1,\ b_1,\ \ldots$ \\
    \bottomrule
  \end{tabular}
\end{table}

The iso-singlet state ($I = 0$), the $\eta$-meson, is described by the
interpolator
\begin{equation}
  O_{\eta}(n) = \frac{1}{\sqrt{2}}\left[ \bar u(n)\, \gamma_5\, u(n)
    + \bar d(n)\, \gamma_5\, d(n) \right].
\end{equation}
The properties of $O_{\pi^0}$ and $O_\eta$ under the parity transformation are
the same as for the interpolators $O_{\pi^\pm}$, i.e., they have $P = -1$.
Concerning charge conjugation, $O_{\pi^0}$ and $O_\eta$ are eigenstates with
$C = +1$. Other mesons differ from the examples we have discussed so far by
their flavor content as well as by spin and parity. We can, e.g., obtain the
interpolator for the strange $K^+$ meson from the interpolator for $\pi^+$, by
replacing the $d$ quark with an $s$ quark,
\begin{equation}
  O_{K^+}(n) = \bar s(n)\, \gamma_5\, u(n) .
\end{equation}
Different spin and parity correspond to different gamma matrices in the
fermionic bilinears. For example we obtain from the interpolator for $\pi^+$
the interpolator for the $\rho^+$ vector meson ($I = 1$, $I_z = +1$, $Q = +e$,
$J = 1$, $P = -1$) by replacing $\gamma_5$ with $\gamma_i$, $i = 1, 2, 3$:
\begin{equation}
  O_{\rho^+}(n)_i = \bar d(n)\, \gamma_i\, u(n) , \qquad i = 1, 2, 3 .
\end{equation}
A general local meson interpolator has the form
\begin{equation}
  O_M(n) = \bar\psi^{(f_1)}(n)\, \Gamma\, \psi^{(f_2)}(n) ,
\end{equation}
where $\Gamma$ is a monomial of gamma matrices, and we have for this equation
switched back to the notation with upper flavor indices $f_i$. In the case of
degenerate flavors ($f_1 = f_2$) combinations of (6.9) are formed to obtain the
desired flavor symmetries, as was done in our interpolators for $\pi^0$ and
$\eta$ in (6.5) and (6.6). In Table~6.1 we list the matrices $\Gamma$ for the
most commonly used interpolators together with the corresponding quantum
numbers. We remark that in addition to the local interpolators gauge-invariant
combinations of quarks and antiquarks at different lattice sites may also be
used. (See Sects.~6.2.2 and 6.2.3.)

\subsection{Meson correlators}

So far we have constructed the interpolator $O_M$ that corresponds to the
operator $\hat{O}_M$ acting in the physical Hilbert space. For the Euclidean
correlator we also need to find the interpolator corresponding to
$\hat{O}_M^\dagger$ which generates the meson state from the vacuum. In order
to identify this operator we formally conjugate the interpolator (6.9) (we drop
the space--time argument $n$ here),
\begin{equation}
  \left( \bar\psi^{(f_1)} \Gamma \psi^{(f_2)} \right)^\dagger
  = -\bar\psi^{(f_2)}\, \bar{\Gamma}^\dagger\, \psi^{(f_1)}
  \equiv -\bar\psi^{(f_2)}\, \gamma_4 \Gamma^\dagger \gamma_4\, \psi^{(f_1)}
  = \pm \bar\psi^{(f_2)}\, \Gamma\, \psi^{(f_1)} .
\end{equation}
The minus sign in the first step comes from the interchange of the Grassmann
variables. In the second step we have used the relation $\bar\psi =
\psi^\dagger\gamma_4$ which holds for the operators in Hilbert space and
performed the equivalent replacement for the Grassmann spinors in our
interpolator. The last step in (6.10) reflects the simple algebraic property
$\gamma_4\Gamma^\dagger\gamma_4 = \pm\Gamma$. Thus, up to a possible overall
sign, the interpolator that corresponds to $\hat{O}_M^\dagger$ is obtained by
interchanging $\psi$'s and $\bar\psi$'s and ordering the $\bar\psi$ to the
left. The sign is irrelevant for the exponential decay in (6.1). Thus, for
mesons we use interpolators
\begin{equation}
  \begin{gathered}
    O_M(n) = \bar\psi^{(f_1)}(n)\, \Gamma\, \psi^{(f_2)}(n) , \\
    \bar{O}_M(m) = \bar\psi^{(f_2)}(m)\, \Gamma\, \psi^{(f_1)}(m) ,
  \end{gathered}
\end{equation}
combined in correlators $\big\langle O_M(n)\, \bar{O}_M(m)\big\rangle$. At the
moment the space--time arguments $m$, $n$ are arbitrary and only later we will
set the time arguments to $m_4 = 0$ and $n_4 = n_t$ to match the form of the
hadron correlator (6.1). In Sect.~6.1.4 we will discuss what to do with the
spatial components $\mathbf{m}$, $\mathbf{n}$.

For evaluating the correlators we have to compute the Grassmann integrals that
appear when calculating the fermionic part $\langle\ldots\rangle_F$ of the
expectation value (compare (5.2)). In this step there is an important
difference between the correlators for iso-triplet operators such as (6.2) and
(6.5) and correlators for an iso-singlet operator such as (6.6). For an
iso-triplet operator of the form $O_T = \bar d\,\Gamma\,u$ we find (summation
convention is used for all indices)
\begin{equation}
  \begin{split}
    \big\langle O_T(n)\, \bar{O}_T(m) \big\rangle_F
    &= \big\langle \bar d(n)\Gamma u(n)\ \bar u(m)\Gamma d(m) \big\rangle_F \\
    &= \Gamma_{\alpha_1\beta_1}\Gamma_{\alpha_2\beta_2}\,
       \big\langle \bar d(n)^{\alpha_1}_{c_1} u(n)^{\beta_1}_{c_1}\,
       \bar u(m)^{\alpha_2}_{c_2} d(m)^{\beta_2}_{c_2} \big\rangle_F \\
    &= -\Gamma_{\alpha_1\beta_1}\Gamma_{\alpha_2\beta_2}\,
       \big\langle u(n)^{\beta_1}_{c_1}\, \bar u(m)^{\alpha_2}_{c_2} \big\rangle_u\,
       \big\langle d(m)^{\beta_2}_{c_2}\, \bar d(n)^{\alpha_1}_{c_1} \big\rangle_d \\
    &= -\Gamma_{\alpha_1\beta_1}\Gamma_{\alpha_2\beta_2}\,
       \big[D_u^{-1}(n|m)\big]^{\beta_1\alpha_2}_{c_1c_2}\,
       \big[D_d^{-1}(m|n)\big]^{\beta_2\alpha_1}_{c_2c_1} \\
    &= - \tr\!\left[ \Gamma\, D_u^{-1}(n|m)\, \Gamma\, D_d^{-1}(m|n) \right].
  \end{split}
\end{equation}
In the third line of this equation we have reordered the Grassmann variables
and then used the fact that the fermionic expectation value factorizes with
respect to the flavors: $\langle\ldots\rangle_F = \langle\ldots\rangle_u\,
\langle\ldots\rangle_d$. Subsequently we applied Wick's theorem (5.36) for
each of the two flavors (compare also (5.54)). This step is often referred to
as fermion contraction.

\begin{figure}[htbp]
  \centering
  \includegraphics[width=0.9\textwidth]{fig61.png}
  \caption*{Fig. 6.1. \emph{Connected (left-hand side plot) and disconnected
    (right-hand side plot) pieces of a meson correlator}.}
  \label{fig:61}
\end{figure}

The Dirac operators $D_u$, $D_d$ for $u$ and $d$ quark differ only by the value
of the mass parameter (compare (5.51)). Often the small difference between the
$u$ and the $d$ quark masses is ignored and one uses $D_u = D_d$, i.e., exact
isospin symmetry. It is, however, important to keep in mind that also in this
case, only Grassmann variables with equal flavor can be contracted with each
other.

The result in the last line of (6.12) has a simple interpretation: The
propagator $D_u^{-1}(n|m)$ propagates a $u$ quark from space--time point $m$ to
the point $n$, while the propagator $D_d^{-1}(m|n)$ transports a $d$ quark in
the opposite direction. Such a contribution is referred to as connected piece
and is depicted in the left-hand side plot of Fig.~6.1. We remark that each of
the individual lines in this figure symbolizes a collection of fermion lines
(cf.\ Fig.~5.1).

In the correlator of an iso-singlet operator $O_S = (\bar u\Gamma u +
\bar d\Gamma d)/\sqrt{2}$, such as (6.6), also another type of contribution
appears. The fermion contractions for this correlator are obtained by following
the same steps as in (6.12),
\begin{multline}
  \big\langle O_S(n)\, \bar{O}_S(m) \big\rangle_F
  = -\frac{1}{2}\tr\!\left[ \Gamma D_u^{-1}(n|m)\Gamma D_u^{-1}(m|n) \right] \\
  + \frac{1}{2}\tr\!\left[ \Gamma D_u^{-1}(n|n) \right]
    \tr\!\left[ \Gamma D_u^{-1}(m|m) \right] \\
  + \frac{1}{2}\tr\!\left[ \Gamma D_u^{-1}(n|n) \right]
    \tr\!\left[ \Gamma D_d^{-1}(m|m) \right] \\
  + u \leftrightarrow d .
\end{multline}
The first type of contribution are the connected pieces we have already
discussed. However, one also gets propagators $D_u^{-1}(n|n)$,
$D_u^{-1}(m|m)$ which transport a $u$ quark from a space time-point back to the
same point. Such terms are called disconnected pieces and are depicted in the
right-hand side plot of Fig.~6.1. Numerically these contributions need more
computational effort and higher statistics than the connected parts and many
studies avoid considering such mesons or drop the disconnected pieces.

We remark that the interpolator $O_{T,I_z=0} = (\bar u\Gamma u - \bar d\Gamma
d)/\sqrt{2}$ for the $I_z = 0$ component of the iso-triplet differs from the
singlet interpolator only by a relative minus sign between the $u$ and the $d$
terms (compare (6.5) and (6.6)). The corresponding correlator is like in
(6.13), but with a minus sign in the third term. In the case of exact isospin
symmetry, $D_u = D_d$, the disconnected pieces cancel. The resulting correlator
is the same as for the other members of the triplet (compare (6.12)), implying
that the masses of all triplet states are degenerate. In nature this
degeneracy is quite accurate, since $m_{\pi^\pm} = 140$\,MeV and $m_{\pi^0} =
135$\,MeV. The small discrepancy is due to slightly different masses for $u$
and $d$ quarks and electroweak corrections.

\subsection{Interpolators and correlators for baryons}

Baryons are objects made out of three valence quarks. Baryon interpolators $O$
with well-defined quantum numbers have to be constructed such that the
corresponding Hilbert space operator $\hat{O}$ in (6.1) projects onto the state
we are interested in. As for the case of mesons we first discuss an example in
detail, the construction of nucleon interpolators, and then address other
baryons. A review on baryon interpolators, going beyond our introductory
presentation, can be found in [2].

The proton $p$ and the neutron $n$ are the $I_z = +1/2$ and $I_z = -1/2$
components of an iso-doublet ($I = 1/2$). Their almost degenerate masses of
$m_p = 938$\,MeV and $m_n = 940$\,MeV demonstrate again that isospin is a good
symmetry, as we already stressed for the pion system. The proton has electric
charge $Q = e$, while the neutron has vanishing charge. Consequently the proton
is a $uud$ state and the neutron is of the type $ddu$. Due to the facts that
the electric charge is irrelevant in QCD and $p$ and $n$ are strongly related
by isospin symmetry, usually one does not distinguish between the two and
refers to them simply as nucleons. We discuss only the $uud$-type interpolator
($I_z = 1/2$). The $I_z = -1/2$ component is obtained by interchanging $u$ and
$d$.

The simplest interpolator for the nucleon $N$ is given by
\begin{equation}
  O_N(n) = \varepsilon_{abc}\, u(n)_a\, u(n)_b^T\, C\, \gamma_5\, d(n)_c .
\end{equation}
In this equation we show explicitly only the color indices $a, b, c$, while for
the Dirac indices we use vector/matrix notation. The transposition $T$ acts on
the Dirac indices and turns the column 4-spinor $u(n)_b$ into a row spinor
$u(n)_b^T$. Summing the color indices with the epsilon tensor makes the
interpolator a color singlet and gauge-invariant, as follows from the action of
a gauge transformation on the fermions (2.30) and (3.31).

The term in the parentheses combines a $u$ and a $d$ quark into a so-called
diquark, using the charge conjugation matrix $C$ and $\gamma_5$. The diquark
has isospin $I = 0$ and spin $J = 0$. The notion diquark here serves to discuss
the quantum numbers and no dynamical meaning is implied. Thus the full
interpolator $O_N$ has $I = 1/2$, $I_z = +1/2$ and $J = 1/2$ as needed for the
nucleon.

The last quantum number we have to discuss is parity, which is $P = +1$ for
proton and neutron. Under the parity transformation $P$ (see (5.72)) our
nucleon interpolator $O_N(n)$ transforms as
\begin{multline}
  O_N(n, n_4) \xrightarrow{\;P\;} \varepsilon_{abc}\,
    \gamma_4 u(-n, n_4)_a\, u(-n, n_4)_b^T\, \gamma_4^T C\gamma_5\gamma_4\,
    d(-n, n_4)_c \\
  = \varepsilon_{abc}\, \gamma_4 u(-n, n_4)_a\, u(-n, n_4)_b^T\, C\gamma_5\,
    d(-n, n_4)_c \\
  = \gamma_4\, O_N(-n, n_4) ,
\end{multline}
where in the second step we used $\gamma_\mu^T C = -C\gamma_\mu$ (compare
(A.23)). As for mesons, the change of the spatial vector $n$ into $-n$ is
irrelevant when we project to zero momentum, i.e., when we sum over all $n$
(see Sect.~6.1.4). The nontrivial transformation of the spinor indices with
$\gamma_4$ can be taken into account by considering the combinations
\begin{equation}
  O_{N^\pm}(n) = \half\left[ O_N(n) \pm O_N^P(n) \right]
  = \varepsilon_{abc}\, P_\pm\, u(n)_a\, u(n)_b^T\, C\gamma_5\, d(n)_c ,
\end{equation}
where the parity projectors $P_\pm$ are defined as
\begin{equation}
  P_\pm = \half (1 \pm \gamma_4) .
\end{equation}
After projection to zero momentum, the interpolators in (6.16) have definite
parity $P = \pm 1$. The interpolator $O_{N^+}$ describes the positive parity
nucleon, while $O_{N^-}$ couples to its negative parity partner $N(1535)$ with
a mass of 1535\,MeV.

Similar to the steps performed for the meson correlator in (6.10), one finds
that the interpolator for the corresponding creation operator is given by (up
to an overall sign)
\begin{equation}
  \bar{O}_{N^\pm}(n) = \varepsilon_{abc}\, \bar u(n)_a\, C\gamma_5\,
    \bar d(n)_b^T\, \bar u(n)_c\, P_\pm .
\end{equation}
Interpolators for other spin-$1/2$ baryons are obtained by changing the flavor
content. Furthermore, different diquark structures can be considered and we
obtain more general spin-$1/2$ interpolators:
\begin{equation}
  \begin{gathered}
    O_{N\pm} = \varepsilon_{abc}\, P_\pm \Gamma_A\, u_a\, u_b^T\, \Gamma_B\, d_c , \\
    \bar{O}_{N\pm} = \varepsilon_{abc}\, \bar u_a\, \Gamma_B\, \bar d_b^T\,
      \bar u_c\, \Gamma_A\, P_\pm , \\
    O_{\Sigma\pm} = \varepsilon_{abc}\, P_\pm \Gamma_A\, u_a\, u_b^T\,
      \Gamma_B\, s_c , \\
    \bar{O}_{\Sigma\pm} = \varepsilon_{abc}\, \bar u_a\, \Gamma_B\, \bar s_b^T\,
      \bar u_c\, \Gamma_A\, P_\pm , \\
    O_{\Xi\pm} = \varepsilon_{abc}\, P_\pm \Gamma_A\, s_a\, s_b^T\, \Gamma_B\,
      u_c , \\
    \bar{O}_{\Xi\pm} = \varepsilon_{abc}\, \bar s_a\, \Gamma_B\, \bar u_b^T\,
      \bar s_c\, \Gamma_A\, P_\pm , \\
    O_{\Lambda\pm} = \varepsilon_{abc}\, P_\pm \Gamma_A\,
      \bigl[ 2\, s_a\, u_b^T\, \Gamma_B\, d_c + d_a\, u_b^T\, \Gamma_B\, s_c
        - u_a\, d_b^T\, \Gamma_B\, s_c \bigr] , \\
    \bar{O}_{\Lambda\pm} = \varepsilon_{abc}\,
      \bigl[ 2\, \bar u_a^T\, \Gamma_B\, \bar d_b\, \bar s_c
        + \bar u_a^T\, \Gamma_B\, \bar s_b\, \bar d_c
        - \bar d_a\, \Gamma_B\, \bar s_b\, \bar u_c \bigr]\,
      \Gamma_A\, P_\pm ,
  \end{gathered}
\end{equation}
where we list the interpolators for nucleon, for the strange $\Sigma$ ($I = 1$,
$S = -1$), $\Xi$ ($I = 1/2$, $S = -2$), and $\Lambda$ ($I = 0$, $S = -1$)
baryons of both parities. Three possible choices for the matrices $\Gamma_A$
and $\Gamma_B$, all giving rise to $J^P = 1/2^+$, are $(\Gamma_A, \Gamma_B) =
(1, C\gamma_5)$, $(\gamma_5, C)$, or $(1, i\gamma_4 C\gamma_5)$.

Baryons with spin $J = 3/2$ can be obtained by using $(\Gamma_A, \Gamma_B) =
(1, C\gamma_j)$. The spatial gamma matrix $\gamma_j$ gives rise to a diquark
with spin $J = 1$ and together with the quark outside the diquark the resulting
interpolator has spin $J = 3/2$ contributions, but also an admixture of spin
$J = 1/2$. For the details of the necessary projection to definite spin we
refer to [2].

Note that all our baryon interpolators $O_{B\pm}$ and $\bar{O}_{B\pm}$ still
have an open Dirac index -- they describe fermions after all. Usually this
index is summed over and one considers correlators of the type ($\alpha$ is
summed)
\begin{equation}
  \big\langle O_{B\pm}(n)_\alpha\, \bar{O}_{B\pm}(m)_\alpha \big\rangle .
\end{equation}
We conclude with displaying as an example the fermion contractions for the
nucleon interpolator (6.14) (we use $P_\pm^2 = P_\pm$ and vector/matrix
notation),
\begin{equation}
  \begin{gathered}
    \big\langle O_{N^\pm}(n)_\alpha\, \bar{O}_{N^\pm}(m)_\alpha \big\rangle_F
    = -\big\langle \bar{O}_{N^\pm}(m)_\alpha\, O_{N^\pm}(n)_\alpha \big\rangle_F \\
    = -\varepsilon_{abc}\varepsilon_{a'b'c'}\,
      \Big\langle \bar u(m)_a C\gamma_5 \bar d(m)_b^T \bar u(m)_c P_\pm \\
    \qquad\qquad u(n)_{c'}\, u(n)_{a'}^T C\gamma_5\, d(n)_{b'} \Big\rangle_F \\
    = \varepsilon_{abc}\varepsilon_{a'b'c'}\, (C\gamma_5)_{\alpha'\beta'}
      (C\gamma_5)_{\alpha\beta} (P_\pm)_{\gamma\gamma'}\,
      D_d^{-1}(n|m)_{\beta'\beta} \times \\
    \qquad \Big[ D_u^{-1}(n|m)_{\alpha'\alpha} D_u^{-1}(n|m)_{\gamma'\gamma}
      - D_u^{-1}(n|m)_{\alpha'\gamma} D_u^{-1}(n|m)_{\gamma'\alpha} \Big] ,
  \end{gathered}
\end{equation}
demonstrating that for baryons only connected pieces occur.

\subsection{Momentum projection}

There is a final step to be implemented for our hadron interpolators: We want
our hadron states to be states of definite spatial momentum $\mathbf{p}$. Thus
we define
\begin{equation}
  \widetilde{O}(\mathbf{p}, n_t) = \frac{1}{|\Lambda_3|}
    \sum_{n\in\Lambda_3} O(n, n_t)\, e^{-\imath a\, n\cdot p} .
\end{equation}
Note that this Fourier transformation runs only over the spatial components
$n$ in the spatial lattice $\Lambda_3 = \{ n = (n_1, n_2, n_3) \mid n_i = 0, 1,
\ldots, N-1 \}$. Accordingly, the momenta $\mathbf{p}$ are spatial momenta with
components $p_i = 2\pi k_i/(aN)$, $k_i = -N/2 + 1, \ldots, N/2$. Thus, the
interpolator $\widetilde{O}(\mathbf{p}, n_t)$ is projected to definite spatial
momentum and is located on a single time slice $n_t$.

It is sufficient to project only one of the two interpolators of a correlation
function to definite momentum, typically the operator $O(n, n_t)$ at the sink.
The source operator $O(0, 0)$ can remain in real space and usually is placed at
the origin. That this is sufficient can be seen by writing the real space
interpolator $O$ as a sum of its Fourier components (compare (A.37)) and using
the fact that states with different momenta are orthogonal to each other such
that only the zero momentum term survives. Note, however, that this is strictly
correct only for the exact expectation values, whereas it is only approximate
for a sum over a finite number of configurations.

Thus, according to (1.21), our final formula for Euclidean hadron correlators
with definite momentum $\mathbf{p}$ reads
\begin{equation}
  \begin{split}
    \big\langle \widetilde{O}(\mathbf{p}, n_t)\, O(0, 0) \big\rangle
    &= \frac{1}{|\Lambda_3|}\sum_{n\in\Lambda_3} e^{-\imath a\, n\cdot p}\,
       \big\langle O(n, n_t)\, O(0, 0) \big\rangle \\
    &= A\, e^{-a n_t E(\mathbf{p})}\left( 1 + \order{e^{-a n_t \Delta E}}
       \right),
  \end{split}
\end{equation}
where the energy $E(\mathbf{p})$ is related to the hadron mass $m_H$ through
the relativistic dispersion relation (in our units the speed of light is
$c = 1$)
\begin{equation}
  E(\mathbf{p}) = \sqrt{m_H^2 + \mathbf{p}^2}\left( 1 + \order{ap} \right),
\end{equation}
and for zero momentum we find $E(0) = m_H$.

Determining the hadron propagators for non-zero momenta allows one to study the
spectral relation between energy and momentum as realized on the lattice. Since
the asymptotic states describe free particles in the lattice world, the
adequate spectral relation is that of free lattice particles. It may be
obtained, e.g., from the free boson (or fermion) lattice propagator in momentum
space. For the simplest, nearest neighbor discretization one identifies the
pole at momentum $(iE(\mathbf{p}), \mathbf{p})$ and finds the relation
($m_H = E(0)$)
\begin{equation}
  \cosh\bigl(a E(\mathbf{p})\bigr) = \cosh(a m_H)
    + \sum_{k=1}^{3} \bigl( 1 - \cos(a p_k) \bigr) ,
\end{equation}
which approaches the continuum relation for vanishing lattice spacing $a$, cf.\
(6.24). Careful analysis of hadron propagators along these lines allows one to
check for the recovery of full Euclidean invariance.

\subsection{Final formula for hadron correlators}

Central for obtaining results to be compared with experiment are the
correlation functions of products of fields. For any product of operators,
denoted by $A$, one has to evaluate
\begin{multline}
  \langle A\rangle = \big\langle \langle A\rangle_F \big\rangle_G
  = \frac{1}{Z}\int D[U]\, e^{-S_G[U]} \\
  \times \int D[\bar\psi,\psi]\, e^{-S_F[\bar\psi,\psi,U]}\,
    A[\bar\psi,\psi,U] .
\end{multline}
The integral over the Grassmann variables can be computed in closed form (see
Chap.~5) and the integral (5.4) over the gauge field remains.

For definiteness let us discuss the example of a propagator for a meson built
out of two quark flavors as in (6.12). After the Grassmann integration we end
up with a ratio of purely bosonic integrals over the gauge fields,
\begin{equation}
  \begin{gathered}
    \big\langle O_T(n)\, \bar{O}_T(m) \big\rangle
    = -\frac{1}{Z}\int D[U]\, e^{-S_G[U]}\, \det[D_u]\det[D_d] \\
    \qquad \times \tr\!\left[ \Gamma D_u^{-1}(n|m)\, \Gamma D_d^{-1}(m|n) \right], \\
    Z = \int D[U]\, e^{-S_G[U]}\, \det[D_u]\det[D_d] .
  \end{gathered}
\end{equation}
We stress again that the fermion determinants depend on $U$ and one has to
perform a Monte Carlo simulation using $Z^{-1}\exp(-S_G[U])\det[D_u]\det[D_d]$
as the distribution weight for the gauge fields. This seems to be a
straightforward task at first sight. However, computing the determinant of the
Dirac operator matrix $D$ is highly nontrivial. Remember that this matrix has
$12|\Lambda|$ rows and columns, where $|\Lambda|$ is the number of lattice
sites. This dimension therefore is typically of the order of millions and a
brute force calculation of the determinant is prohibitively expensive.

\subsection{The quenched approximation}

Following the first applications of the Monte Carlo method to lattice QCD, most
of the results in the 1980s and 1990s were obtained in the quenched
approximation.\footnote{The terms come from quenching -- rapid cooling -- of
steel, where the carbon atoms are frozen to their random positions.} In
particular the successful determination of the ground state spectrum of hadrons
with light quarks is a milestone of this approach.

Before discussing the quenched approximation, let us briefly remind ourselves
of the physical interpretation of the determinant factor. From the hopping
expansion discussed in Chap.~5, we found that it may be written as an effective
action, which is a sum over closed loops of gauge field link variables. In
essence it describes the fermionic vacuum where virtual pairs of quarks and
antiquarks are created and annihilated. These quarks are usually called sea
quarks, since they are related to Dirac's picture of a fermion sea, where the
holes denote anti-fermions. Of course sea quarks are identical to generic
quarks and we only call them such for clarification of the subsequent
approximation.

Putting the fermion determinants to unity by hand corresponds to neglecting the
effect of the vacuum quark loops (Sect.~5.3.2). Actually, since in the expansion
of the determinant the leading term of the effective action is a sum over
plaquettes, part of the effective action is intrinsically included in the gauge
action by redefining the gauge coupling. This is no problem, since the gauge
coupling itself is a bare parameter and only physical quantities are of
relevance.

Omitting the determinants defines the quenched approximation. Instead of using
the full expression for, e.g., the meson correlator (6.27), one computes
\begin{equation}
  \begin{gathered}
    \big\langle O_T(n)\, \bar{O}_T(m) \big\rangle_{\mathrm{quenched}}
    = -\frac{1}{Z}\int D[U]\, e^{-S_G[U]}\,
      \tr\!\left[ \Gamma D_u^{-1}(n|m)\, \Gamma D_d^{-1}(m|n) \right], \\
    Z = \int D[U]\, e^{-S_G[U]} .
  \end{gathered}
\end{equation}
In this approximation one constructs the Markov chain of gauge configurations
as for pure gauge theory (compare Chap.~3) and evaluates the quark propagators
on those configurations. Combinations of the quark propagators build up the
hadron correlators as discussed in the last section. In order to distinguish
these propagating quarks from the sea quarks, and since they are responsible
for the quantum numbers of the hadron, they are called valence quarks. The
quenched approximation is therefore also called valence approximation.

The quenched approximation may be considered as the limit, in which the sea
quarks become infinitely heavy and therefore cannot be generated from the
vacuum as particle--antiparticle pairs. This can be seen from the hopping
expansion for the fermion determinant (5.65). Since the expansion parameter
$\kappa$ defined in (5.55) vanishes for $m \to \infty$, the determinant becomes
1 in this limit.

\begin{figure}[htbp]
  \centering
  \includegraphics[width=0.9\textwidth]{fig62.png}
  \caption*{Fig. 6.2. \emph{Sample of quark lines contributing in hadron
    propagation (mesons l.h.s., baryons r.h.s.). In the quenched approximation
    only contributions in the upper two rows, i.e., without sea quark loops
    ($=$ closed loops) contribute. The diagrams in the second row appear in the
    quenched approximation and resemble the dynamical diagrams in the bottom
    rows. They are called hairpin diagrams and give rise to logarithmic
    singularities in the quenched approximation}.}
  \label{fig:62}
\end{figure}

A simulation including the determinant and therefore allowing for the full
dynamical vacuum structure of fermions is called a simulation with dynamical
quarks, in contrast to the quenched simulation. Figure 6.2 gives examples for
contributing quark lines in the quenched and the fully dynamical simulation. In
full QCD there is no difference between sea quarks and valence quarks. However,
in our computer laboratory we may give them different masses. Such intermediate
stages, where the masses of the sea quarks and the valence quarks differ, are
called partial quenching, and are sometimes useful for a better control of the
various systematic uncertainties.

We remark that neither the quenched nor the partially quenched approximation
are proper quantum field theories with a valid Hilbert space construction and
correct positivity properties for the fermions. This leads to some problems in
analyzing in particular the results toward small quark masses. Forbidding
dynamical fermion loops also prevents one from studying resonances. However,
these provide a conceptual problem anyhow. They are not asymptotic states of
the field theory and one has to introduce intricate tools for their
identification and the determination of scattering amplitudes in general
(cf.\ Chap.~11). We expect, however, that resonances should in the quenched
case to leading order be observed as bound states at nearby energy values.

Since dynamical simulations will be treated in Chap.~8, in this chapter we will
show exclusively results for the hadron spectrum in the quenched situation.
Obviously these are of much higher statistical precision than those obtained
from simulations with dynamical quarks, albeit with the intrinsic uncertainty
of the effect of quenching. Special features due to quenching (like occurrence
of extra singularities and ghosts) will be discussed in Sect.~6.4.3.

We stress, however, that most of the formulas and techniques presented in this
chapter apply to both quenched and dynamical simulations.

\section{Strategy of the calculation}

In this section we discuss techniques for computing and storing the quark
propagators needed for the hadron correlators presented in the last section. We
describe the construction of realistic quark sources which improve the overlap
of the lattice hadron interpolators with the physical states. The problem of
exceptional configurations and its amelioration through smoothing of the gauge
field are addressed.

\subsection{The need for quark sources}

A complete quark propagator is a matrix as large as the fermion matrix itself.
Depending on the lattice action, the Dirac operator may be sparse, with many
vanishing entries, but the propagator $D^{-1}$ is not and consists of
$\order{10^{12}}$ complex numbers. Each entry $D^{-1}(n|m)^{\beta\alpha}_{ba}$
connects a source point $(m, \alpha, a)$ with a sink point $(n, \beta, b)$.
Here the notion ``point'' relates to site- as well as Dirac and color indices.
Since the matrix provides the information for propagation in one particular
gauge field configuration, its entries will be highly correlated and this
correlation reduces the information content. So for two reasons one does not
try to store the complete propagator matrix: It is wasteful to calculate all
entries and it would need too much computer memory.

However, if we consider instead just the propagator from a fixed site $m_0$, a
fixed Dirac index $\alpha_0$, and a fixed color index $a_0$, to any site of the
lattice, this is just one column of the inverse Dirac operator,
\begin{equation}
  D^{-1}(n|m_0)^{\beta\alpha_0}_{b a_0}
  = \sum_{m,\alpha,a} D^{-1}(n|m)^{\beta\alpha}_{ba}\,
    S_0^{(m_0,\alpha_0,a_0)}(m)^{\alpha}_{a} ,
\end{equation}
where we have introduced the so-called point sources
\begin{equation}
  S_0^{(m_0,\alpha_0,a_0)}(m)^{\alpha}_{a}
  = \delta(m - m_0)\, \delta_{\alpha\alpha_0}\, \delta_{aa_0} .
\end{equation}
The expression (6.29) must be evaluated for 12 sources, one for each
combination $\alpha_0$, $a_0$ of Dirac and color indices. The point sources are
placed at the position $m_0$ of the source interpolator $O_H(m_0)$ and the
column $D^{-1}(n|m_0)$ is the propagator to the position $n$ of the sink
operator $\bar{O}_H(n)$.

Since mesons contain a backward running antiquark (compare Fig.~6.1), one also
needs the propagator in the other direction, i.e., from $n$ to $m_0$. Naively
one would expect that this requires placing sources also at $n$. However, we
can utilize $\gamma_5$-hermiticity, as discussed in Sect.~5.4.3, in order to
obtain the backward running propagator for free.

Using this symmetry property we may write (repeated indices are summed)
\begin{equation}
  \gamma_5 D^{-1} \gamma_5 = D^{-1\,\dagger}
  \quad\Longleftrightarrow\quad
  (\gamma_5)_{\alpha\alpha'}\, D^{-1}(m_0|n)^{\alpha'\beta'}_{dc}\,
  (\gamma_5)_{\beta'\beta} = D^{-1}(n|m_0)^{\beta\alpha\,*}_{cd} ,
\end{equation}
and so obtain the backward running propagator. This can, e.g., be used to write
the iso-triplet pion propagator (6.12) as
\begin{multline}
  \big\langle O_T(n)\, \bar{O}_T(m_0) \big\rangle_F
  = -\tr\!\left[ D^{-1}(n|m_0)\, \gamma_5\, D^{-1}(m_0|n)\, \gamma_5 \right] \\
  = -\sum_{\alpha,\beta,c,d} \left| D^{-1}(n|m_0)^{\alpha\beta}_{cd}
    \right|^2 .
\end{multline}
This example shows explicitly that for the pion, and similarly for other
mesons, only the forward propagator $D^{-1}(n|m_0)$ needs to be calculated. For
baryons we only need forward propagating quarks anyhow (compare (6.21)).

\subsection{Point source or extended source?}

In order to get clear and strong correlation signals, which allow a reliable
analysis, we have to optimize the interpolating fields. Although any operator
with the correct quantum numbers contributes to the physical state, some may be
more important than others. The overlap can be improved considerably by
providing more realistic spatial wave functions. We present the idea using a
meson interpolator as an example, but the generalization to baryons is
straightforward.

We consider a meson interpolator on a single fixed time slice $n_t$, located at
$\mathbf{n}_0 = (n_0, n_t)$. A more general wave function may be introduced by
writing
\begin{equation}
  O_M(\mathbf{n}_0) = \sum_{n_1,n_2}
    F(\mathbf{n}_0; n_1, n_2)^{\alpha_1\alpha_2}_{a_1a_2}\,
    \psi^{(f)}(n_1, n_t)^{\alpha_1}_{a_1}\,
    \bar\psi^{(f)}(n_2, n_t)^{\alpha_2}_{a_2} ,
\end{equation}
where we have introduced a distribution function $F(\mathbf{n}_0; n_1, n_2)$
which combines field values $\psi(n_1, n_t)$ and $\bar\psi(n_2, n_t)$ at
spatial positions $n_1$, $n_2$ in the vicinity of $\mathbf{n}_0$. In the
simplest case, where the complete interpolator is localized on the single site
$\mathbf{n}_0 = (n_0, n_t)$, $F$ reduces to
\begin{equation}
  \begin{gathered}
    F(\mathbf{n}_0; n_1, n_2)^{\alpha_1\alpha_2}_{a_1a_2}
    = \delta(n_0 - n_1)\, \delta_{\alpha_0\alpha_1}\, \delta_{a_0a_1}\,
      \Gamma_{\alpha_0\beta_0}\, \delta(n_0 - n_2)\, \delta_{\beta_0\alpha_2}\,
      \delta_{a_0a_2} \\[2mm]
    \Longrightarrow\qquad
    O_M(\mathbf{n}_0) = \psi^{(f)}(n_0, n_t)^{\alpha_0}_{a_0}\,
      \Gamma_{\alpha_0\beta_0}\, \bar\psi^{(f)}(n_0, n_t)^{\beta_0}_{a_0} ,
  \end{gathered}
\end{equation}
where $\Gamma$ is an element of the Clifford algebra.

However, a more realistic wave function can be obtained by choosing a less
trivial function $F$. In order to be able to work with only one set of 12
sources one usually uses factorizable functions ($\alpha_0$, $\beta_0$, $a_0$
are summed)
\begin{equation}
  F(\mathbf{n}_0; n_1, n_2)^{\alpha_1\alpha_2}_{a_1a_2}
  = S_i^{(n_0,\alpha_0,a_0)}(n_1)^{\alpha_1}_{a_1}\,
    \Gamma_{\alpha_0\beta_0}\,
    S_k^{(n_0,\beta_0,a_0)}(n_2)^{*\alpha_2}_{a_2} .
\end{equation}
We have introduced the subscripts $i$ and $k$ to indicate that the source
functions for $\psi$ and $\bar\psi$ may differ. Using the form (6.35) one can
rewrite (6.33) in terms of so-called smeared fermions. To indicate the type of
source we mark the smeared fermions with a lower index $i$ or $k$,
respectively:
\begin{equation}
  \begin{gathered}
    \psi_k^{(f')}(n_0, n_t)^{\alpha_0}_{a_0}
    \equiv \sum_{n_2,a_2} S_k^{(n_0,\alpha_0,a_0)}(n_2)^{*\alpha_2}_{a_2}\,
      \psi^{(f')}(n_2, n_t)^{\alpha_2}_{a_2} , \\
    \psi_i^{(f)}(n_0, n_t)^{\alpha_0}_{a_0}
    \equiv \sum_{n_1,a_1} S_i^{(n_0,\alpha_0,a_0)}(n_1)^{\alpha_1}_{a_1}\,
      \psi^{(f)}(n_1, n_t)^{\alpha_1}_{a_1} ,
  \end{gathered}
\end{equation}
This form makes it obvious that the smeared field at $\mathbf{n}_0$ is combined
from the original fields at positions $n_i$. The meson operator in terms of the
smeared fermions becomes simply (this is (6.34), but now for smeared fields):
\begin{equation}
  O_M(\mathbf{n}_0) = \psi_i^{(f)}(n_0, n_t)^{\alpha_0}_{a_0}\,
    \Gamma_{\alpha_0\beta_0}\,
    \psi_k^{(f')}(n_0, n_t)^{\beta_0}_{a_0} .
\end{equation}
Introducing these smearing functions for sources $S_i$, located at $m_0$, and
sinks $S_k$, located at $n_0$, leads to the smeared quark propagator
\begin{equation}
  \begin{split}
    G_{ki}(n_0|m_0)^{\beta_0\alpha_0}_{b_0a_0}
    &\equiv \big\langle \psi_k^{(f)}(n_0)^{\beta_0}_{b_0}\,
       \bar\psi_i^{(f)}(m_0)^{\alpha_0}_{a_0} \big\rangle_F \\
    &= \sum_{m,n} S_k^{(n_0,\beta_0,b_0)}(n)^{*\beta}_{b}\,
       S_i^{(m_0,\alpha_0,a_0)}(m)^{\alpha}_{a}\,
       \big\langle \psi(n)^{\beta}_{b}\, \bar\psi(m)^{\alpha}_{a} \big\rangle_F \\
    &= \sum_{n,m} S_k^{(n_0,\beta_0,b_0)}(n)^{*\beta}_{b}\,
       \big[D_f^{-1}(n|m)\big]^{\beta\alpha}_{ba}\,
       S_i^{(m_0,\alpha_0,a_0)}(m)^{\alpha}_{a} .
  \end{split}
\end{equation}
In the construction of the meson and baryon correlators, as discussed in
Sect.~6.1, one now simply replaces the original quark propagator $D_f^{-1}$ by
the smeared propagator $G_{ki}$. In an actual numerical calculation one
computes the propagator from a smeared source $S_i$ to a point-like sink
according to
\begin{equation}
  \sum_{m,\alpha,a} D^{-1}(n|m)^{\beta\alpha}_{ba}\,
    S_i^{(m_0,\alpha_0,a_0)}(m)^{\alpha}_{a} .
\end{equation}
The smearing of the sink can then be done in the subsequent construction of the
hadron propagators by applying the smearing operator presented in the next
subsection.

\subsection{Extended sources}

For smearing functions $S$ that are not gauge-covariant, one has to fix the
gauge in the time slice using the methods from Sect.~3.2. Various source shapes
have been used in that context. The extreme among those is a source constant in
a time slice, the so-called wall-source. Attractive are Gaussian shapes
motivated by intuition. The final incentive always is the optimization of the
signal.

A gauge-covariant source, with a shape similar to a Gaussian, is obtained by
Jacobi smearing [3, 4]. One acts with a smearing operator $M$ on the point
source $S_0$ of (6.30) in time slice $n_t$ to obtain the smeared source:
\begin{equation}
  S^{(n_0,\alpha_0,a_0)} = M\, S_0^{(n_0,\alpha_0,a_0)}, \qquad
  M = \sum_{n=0}^{N} \kappa^n H^n .
\end{equation}
The operator $H$ is essentially the spatial part of the Wilson term (5.50)
without the constant piece ($n_t$ is fixed to the time slice of the source),
\begin{equation}
  H(n, m) = \sum_{j=1}^{3}\left[ U_j(n, n_t)\, \delta(n + \hat{\jmath}, m)
    + U_j(n - \hat{\jmath}, n_t)^\dagger\, \delta(n - \hat{\jmath}, m) \right].
\end{equation}
Note that $H$ and therefore also $M$ are hermitian and act only on the color
indices, but are trivial in Dirac space. The operation connects different sites
of the time slice to the central site with gauge transporters. Jacobi smearing
has two free parameters: the number of smearing steps $N$ and the positive real
parameter $\kappa$. These two parameters can be used to adjust the profile
(width) of the source. Also combining quark sources of different widths has
proved to be efficient for optimizing the hadron propagation signals
(cf.\ Sect.~6.3.3).

\subsection{Calculation of the quark propagator}

Having prepared our sources $S$, we need to compute the propagator $D^{-1}$
acting on the source according to (6.39), i.e., we need $G = D^{-1}S$. Put
differently, one has to solve the system of equations
\begin{equation}
  D\, G = S ,
\end{equation}
where $D$ is the Dirac matrix operator, $S$ the source vector, and $G$ the
unknown propagator vector. This can be done using iterative methods. In order
to get acquainted with that approach let us briefly discuss the grandmother of
those algorithms: As shown in the discussion of the hopping expansion (compare
(5.57)), after rescaling the fermions we may split $D$ into a constant part and
a nontrivial term,
\begin{equation}
  D \equiv 1 - Q .
\end{equation}
Then a series expansion
\begin{equation}
  G = (1 - Q)^{-1} S = (1 + Q + Q^2 + Q^3 + \ldots)\, S
\end{equation}
leads to the iteration prescription
\begin{equation}
  G^{(0)} = S , \qquad G^{(i+1)} = S + Q\, G^{(i)} .
\end{equation}
This is the simple Jacobi iteration which converges when the largest eigenvalue
of $Q$ is smaller than 1 in magnitude. In each iteration step one only has to
apply the matrix $Q$ to the vector $G^{(i)}$ and add $S$ until convergence is
obtained. Relaxation methods like the Gauss--Seidel method interpolate between
subsequent steps and may improve convergence. Both methods belong to the class
of stationary methods which solve for each of the variables locally, by keeping
the current values of the other variables fixed. Other variants derived from
the Gauss--Seidel method, but converging an order of magnitude faster, are the
SOR (successive overrelaxation) and SSOR (symmetric SOR) methods, which
introduce extrapolation parameters [5].

Most calculations, however, use a variant of the conjugate gradient (CG) method
[5, 6]. This is a so-called non-stationary method and the concept is to
iteratively minimize a quadratic functional, equivalent to finding the solution
of the system of equations. For real symmetric, positive definite $N\times N$
matrices $A$ the function
\begin{equation}
  Q(x) = \half\, x^T A\, x - x^T b
\end{equation}
is (up to a constant) a positive definite quadratic form assuming its minimum
at $x^*$, where its gradient vanishes,
\begin{equation}
  \left.\frac{\partial Q(x)}{\partial x}\right|_{x = x^*}
  = A\, x^* - b = 0 .
\end{equation}
Thus $x^*$ is the solution of the system $Ax = b$. The CG method works by
iteratively constructing search direction vectors $p^{(i)}$ and iterates
$x^{(i)}$ such that in each step $Q(x^{(i)} + \alpha_i p^{(i)})$ is minimized as
a function of the real parameter $\alpha_i$, leading to the next iterate
$x^{(i+1)} = x^{(i)} + \alpha_i p^{(i)}$. The vectors $p^{(i)}$ are orthogonal
to $A\,p^{(i-1)}$ (and all previous $A\,p^{(j)}$) sequentially build up a
vector space $K^{(i)} = \mathrm{span}(p^{(0)}, \ldots, p^{(i)})$, which is
called a Krylov subspace. All variants of the CG method are based on such an
implicit construction of an orthogonal basis for the Krylov subspace. Methods
like the Lanczos algorithm for finding the eigensystem of a matrix work with
this concept as well and the corresponding iteration may be related to the CG
iteration [5].

The iterates $x^{(i)}$ minimize $Q$ in the spaces $K^{(i)}$ and approach the
solution vector $x^*$ in at most $N$ steps. However, often much fewer steps
give sufficiently accurate results. This may be checked by computing the norm
$\|A\,x^{(i)} - b\|$, and the process is terminated once the norm is smaller
than the requested accuracy $\varepsilon$. Only scalar products between vectors
and matrix--vector multiplications have to be computed and only vectors have to
be kept in storage. Both are important advantages for the large, but often
sparse, matrices associated with the Dirac operators.

In its original version CG works only for positive definite symmetric matrices
$A$. However, for non-symmetric matrices the residual vectors cannot be made
orthogonal within these limitations. The Dirac operator $D$ is not positive
definite hermitian; it has complex eigenvalues. For non-symmetric, general
matrices a variant called Bi-Conjugate Gradient (Bi-CGR) does the job. There
two sequences of search directions and residual vectors are generated, which
obey a bi-orthogonality relation. Even better, improving convergence
properties, the so-called Bi-CGR Stabilized (Bi-CGStab) method [7] and another
improvement, Bi-CGStab(2) by [8], have been introduced. For details we refer to
the original papers and to [5].

We remark that one could utilize $D^{-1} = D^\dagger(D D^\dagger)^{-1}$ and
apply an algorithm for hermitian matrices to compute $(D D^\dagger)^{-1}$.
Furthermore, for $\gamma_5$ hermitian Dirac operators one could also invert the
hermitian matrix $D\gamma_5$ and use $D^{-1} = \gamma_5 (D\gamma_5)^{-1}$. In
practice, however, it turned out that the strategy of using an algorithm for
hermitian matrices does not necessarily lead to faster convergence.

Here we discuss only the Bi-CGStab algorithm along the lines of [5]. In order
to solve the matrix equation (6.47) for the unknown vector $x^*$, one starts
with a guessed vector $x^{(0)}$. The pseudo-code for the iteration with
preconditioning (see below) then may be written as in Table~6.2. For the
application of Bi-CGStab to our problem (6.42), one sets $A = D$, $b = S$, and
obtains the desired propagator vector $G = D^{-1}S$ as
$G = \lim_{i\to\infty} x^{(i)}$.

\begin{table}[htbp]
  \centering
  \caption*{Table 6.2. Pseudocode [5] for the Bi-CGStab algorithm with
    preconditioning matrix $M$. For no preconditioning set $M = 1$ and simplify
    the code}
  \small
  \begin{tabular}{@{}l@{}}
    $r^{(0)} = b - A\, x^{(0)}$ \\
    $\tilde{r} = r^{(0)}$ \quad (for example) \\
    for $i = 1, 2, \ldots$ iterate: \\
    \quad $\rho_{i-1} = \tilde{r}^{\dagger} \cdot r^{(i-1)}$ \\
    \quad if $\rho_{i-1}$ equals 0, the method fails. \\
    \quad if $i = 1$ then \\
    \quad\quad $p^{(1)} = r^{(0)}$ \\
    \quad else \\
    \quad\quad $\beta_{i-1} = \alpha_{i-1}\, \rho_{i-1} / (\omega_{i-1}\,
      \rho_{i-2})$ \\
    \quad\quad $p^{(i)} = r^{(i-1)} + \beta_{i-1}\, (p^{(i-1)} - \omega_{i-1}\,
      v^{(i-1)})$ \\
    \quad end if \\
    \quad solve $M\, \hat{p} = p^{(i)}$ \\
    \quad $v^{(i)} = A\, \hat{p}$ \\
    \quad $\alpha_i = \rho_{i-1} / (\tilde{r}^{\dagger} \cdot v^{(i)})$ \\
    \quad $s = r^{(i-1)} - \alpha_i\, v^{(i)}$ \\
    \quad if the norm of $s$ is small enough, set $x^{(i)} = x^{(i-1)} +
      \alpha_i\, \hat{p}$ and stop. \\
    \quad solve $M\, \hat{s} = s$ \\
    \quad $t = A\, \hat{s}$ \\
    \quad $\omega_i = t^{\dagger} \cdot s / (t^{\dagger} \cdot t)$ \\
    \quad $r^{(i)} = s - \omega_i\, t$ \\
    \quad $x^{(i)} = x^{(i-1)} + \alpha_i\, \hat{p} + \omega_i\, \hat{s}$ \\
    repeat until convergence \ldots
  \end{tabular}
\end{table}

We remark that $\tilde{r}$ remains constant and one may also choose other
$\tilde{r}$, as long as $\rho_0 \neq 0$. For repeating the iteration one needs
$\omega_i \neq 0$ (the Bi-CGStab(2) algorithm helps in this respect [8]). One
stops the iterations as soon as $r^{(i)\dagger} \cdot r^{(i)} < \varepsilon$ for
some requested accuracy $\varepsilon$. The algorithm needs two matrix--vector
multiplications with $A$, two inversions of the preconditioner $M$, and four
inner products for each iteration step.

For the simple CG algorithm convergence is guaranteed for positive definite
symmetric matrices, at least after $N$ steps for matrices of dimension $N$, but
usually significantly faster. For Bi-CGStab convergence may be faster than that
of other CG methods, but not necessarily, and in some cases it even breaks
down. Other starting values or other methods have to be chosen in such
situations.

So-called preconditioning methods sometimes allow one to accelerate the
convergence, depending on the Dirac operator. A preconditioner is a suitable
matrix $M$ which we use to transform the system to
\begin{equation}
  M^{-1} A\, x = M^{-1} b .
\end{equation}
Let us assume that we have a matrix $M$, which is numerically cheap to invert,
i.e., to solve $M\hat{s} = s$ for $\hat{s}$, and which approximates $A$ in some
way. Then the spectral properties of $M^{-1}A$ may be more favorable, allowing
faster convergence of the iterative solution. This may be true in particular,
if the small eigenvalues of $M$ agree with those of $A$. For a more complete
discussion of such methods cf.\ [5, 9].

When inverting the Dirac operators for different quark masses, the matrices
differ just by a constant in the diagonal. One then can simultaneously solve
the set of equations for a set of quark masses with little extra cost. The rate
of convergence is determined by the convergence for the smallest mass. Such a
multi-mass algorithm has been discussed in [10, 11].

Since the determination of the quark propagators is numerically quite
expensive, one often stores these for each gauge configuration sampled and
computes various hadron observables from them in the later analysis.

\subsection{Exceptional configurations}

So far we have not taken into account the Dirac operator matrix $D[U]$ which
depends on the gauge field $U$. This implies that the properties of $D[U]$ will
change as one changes $U$. In particular, certain fluctuations of the gauge
field, so-called exceptional configurations, lead to small eigenvalues of
$D[U]$ which make the numerical inversion of $D[U]$ problematic.

To understand this phenomenon in more detail let us consider the eigenvalues of
$D[U]$. For quarks with bare mass parameter $m$ they are given by
\begin{equation}
  m + \lambda_i[U] ,
\end{equation}
where $\lambda_i[U]$ are the eigenvalues of the massless Dirac operator. In
general these are complex numbers, but also real eigenvalues
$\lambda_i[U] = r_i[U]$ are possible. If such a real eigenvalue becomes
negative and in magnitude similar to $m$, then the Dirac operator has a very
small eigenvalue. In this case the numerical inversion of $D[U]$ breaks down,
and only increasing the quark mass $m$ resolves the problem. Thus the
exceptional configurations limit the quark masses one can reach. For example in
quenched simulations with Wilson fermions it turns out that in a typical
simulation\footnote{Theoretical arguments show that only for $\kappa < 1/8$ one
can be sure that there are no exceptional configurations.} (e.g., for lattice
size $32^3 \times 64$) the quark mass has to be chosen large enough such that
the pion mass is at least 300\,MeV instead of the physical 135\,MeV.

We remark that for chiral fermions, where the Dirac operator obeys the
Ginsparg--Wilson equation, a concept which we discuss in Chap.~7, the real
eigenvalues do not fluctuate and the problem with exceptional configurations
does not occur; therefore realistic pion masses can be reached. Also
introducing a twisted mass term (see Chap.~10) solves the problem with
exceptional configurations, since it expels the spectrum from a strip along the
real axis.

In practical simulations one finds that the fluctuations of the eigenvalues are
coupled to the fluctuations of the gauge field. These fluctuations can be
damped by either using improved gauge field actions (see Chap.~9), larger
lattices, or by smoothing techniques which we discuss in the next subsection.

\subsection{Smoothing of gauge configurations}

When aiming for correlation functions one is mainly interested in the long
distance behavior. On the other hand, typical for gauge theories are violent
short distance fluctuations of the gauge field. One can considerably improve
the correlation signal by smoothing or smearing the gauge field, either only in
the time slices or even in space and time. This also helps with the exceptional
configurations discussed in the preceding subsection.

When smoothing or smearing one typically replaces the link variables by local
averages over short paths connecting the link's endpoints. This is a gauge
covariant procedure and one does not have to fix the gauge. The operators and
propagators are then constructed on the smeared configurations. As long as the
smearing is local enough, i.e., the smearing operator combines only a fixed
number of links, the long distance correlation signals should not be affected
in the continuum limit.

Smearing algorithms are all averaging products of links along certain paths
connecting the endpoints of a given link. For SU(2) this average is
proportional to a group element, but for SU(3) this is not the case. Therefore
one has to project the average to an SU(3) matrix. We briefly mention here only
three variants of such smearing algorithms.

APE smearing [12]: In that case the average is over the original link $U_\mu$
and over the six perpendicular staples (visualize!) connecting its endpoints,
\begin{equation}
  \begin{gathered}
    V_\mu(n) = (1 - \alpha)\, U_\mu + \frac{\alpha}{6}\sum_{\nu\neq\mu}
      C_{\mu\nu}(n) , \\
    C_{\mu\nu}(n) = U_\nu(n) U_\mu(n+\hat\nu) U_\nu(n+\hat\mu)^\dagger
      + U_\nu(n-\hat\nu)^\dagger U_\mu(n-\hat\nu)
      U_\nu(n-\hat\nu+\hat\mu) ,
  \end{gathered}
\end{equation}
where the real parameter $\alpha$ may be adjusted depending on the gauge
coupling. The projection of the sum to SU(3) is usually done by maximizing
$\Re\,\tr[X\, V_\mu(n)^\dagger]$ for $X \in \mathrm{SU}(3)$ and using $X$ as new
link variable $U'_\mu(n)$.

HYP smearing [13]: In that approach the average is over paths lying within the
hypercubes containing the link variable. Again the sum has to be projected to
SU(3).

Stout smearing [14]: This method uses a particular way of projection, defining
the new link after a smearing step by
\begin{equation}
  U'_\mu(n) = e^{\imath Q_\mu(n)}\, U_\mu(n) .
\end{equation}
$Q_\mu(n)$ is a traceless, hermitian matrix constructed from staples,
\begin{equation}
  \begin{gathered}
    Q_\mu(n) = \frac{\imath}{2}\left[ \Omega_\mu(n) - \Omega_\mu(n)^\dagger
      - \frac{1}{3}\tr\!\left( \Omega_\mu(n) - \Omega_\mu(n)^\dagger \right)
      \right], \\
    \Omega_\mu(n) = \left( \sum_{\nu\neq\mu} \rho_{\mu\nu} C_{\mu\nu}(n)
      \right) U_\mu(n)^\dagger .
  \end{gathered}
\end{equation}
The $C_{\mu\nu}$ are the same as in (6.50) and the real weight factors
$\rho_{\mu\nu}$ are tunable parameters. Common choices are taking them
constant, $\rho_{\mu\nu} = \rho$, or smearing only the spatial links:
$\rho_{\mu 4} = \rho_{4\mu} = 0$, $\rho_{nm} = \rho$. The new links have gauge
transformation properties like the original ones. This approach of getting an
SU(3) element from a sum of matrices may be used for other combinations of
paths too, of course. The advantage of stout smearing is that $U'_\mu(n)$ is
differentiable with respect to the link variables. This is a prime benefit in
applications like the hybrid Monte Carlo method for dynamical fermions
discussed in Chap.~8.

Such link smearing methods lead to what is often called a fat link for obvious
reasons. There is a variety of other such suggestions, among them the FLIC link
[15] or differentiable variants of the HYP smearing [16, 17].

All such smearing steps can be iterated. However, the smearing then affects not
just neighboring variables but extends over increasingly larger distances.
Eventually the asymptotic behavior of the propagators will be affected
rendering such results problematic. Anyway, most of the positive effect of
smearing is achieved already after the first smearing step(s).

\section{Extracting hadron masses}

Once the quark propagators are computed as described in the last section, they
can be combined to the hadron correlators discussed in Sect.~6.1. From these
correlators one can extract the corresponding hadron masses and in this section
we present the techniques necessary for this step.

In spectroscopy calculations we analyze hadron correlators of the form (6.1).
In these formulas energies $E$, masses $m$, and momenta $\mathbf{p}$ always
come with a factor of the lattice constant $a$, making the products $aE$,
$am$, and $a\mathbf{p}$ dimensionless. It is convenient to drop the factor of
the lattice constant and use so-called lattice units. In lattice units (6.1)
assumes the form (6.53) below. For going back to physical units, an energy $E$
has to be replaced by $aE$, and equivalently for masses and momenta (see
Sect.~6.4.2). Unless stated otherwise, from now on we use lattice units for
masses and energies.

\subsection{Effective mass curves}

In the analysis of the mass spectrum one studies the correlation function of an
operator $\hat{O}$ with the quantum numbers of a particular hadron or set of
hadrons. However, these operators usually do not create eigenstates of the
Hamiltonian. They have contributions from several eigenstates and the
correlation function is the spectral sum (1.21), discussed in Chap.~1. If the
sink operator is projected to zero momentum (by summation over the respective
time slices) we find
\begin{equation}
  C(n_t) \equiv \big\langle \widetilde{O}(\mathbf{0}, n_t)\, O(0, 0)
    \big\rangle = \sum_k \bra{0}\hat{O}\ket{k}\,\bra{k}\hat{O}^\dagger\ket{0}\,
    e^{-n_t E_k} .
\end{equation}
For finite lattice size the values $E_k$ are discrete. If the operators couple
to single particle intermediate states, the masses $m_k$ of those particles
will correspond to the low-lying energy values, i.e., $E_k = m_k$. If the
operators couple to two- or more-particle intermediate states, the energy
values $E_k$ are related to masses and relative momenta of the involved
particles.

Due to the exponential form of the terms in the sum (6.53), the relative
contribution of higher lying energy states varies with $n_t$. Only for large
$n_t$ the lowest energy state dominates. For smaller $n_t$ the mixing with
other states blurs the picture and the sum (6.53) has many relevant
contributions,
\begin{equation}
  C(n_t) = A_0\, e^{-n_t E_0} + A_1\, e^{-n_t E_1} \ldots .
\end{equation}
In Sect.~6.3.3 we will discuss methods to analyze the case when more than one
energy level has to be determined. For the moment let us consider the simplest
situation where one is interested only in the lowest energy $E_0$ corresponding
to the ground state.

The left-hand side plot of Fig.~6.3 shows an example for the correlation
function of a mesonic operator with the quantum numbers of the pion in a sample
calculation. The linear regions in the log-plot correspond to the range of
$n_t$-values where a single exponential behavior
$\sim \exp(-n_t E_0)$ dominates.

\begin{figure}[htbp]
  \centering
  \includegraphics[width=0.9\textwidth]{fig63.png}
  \caption*{Fig. 6.3. \emph{Result of a Monte Carlo simulation on a
    $16^3 \times 32$ lattice at lattice spacing $a \approx 0.15$~fm. L.h.s.:
    log-plot for the pion correlation function; r.h.s.: effective mass plot (in
    lattice units). The different sets correspond to different values of the
    quark mass in lattice units. The points are connected to guide the eye}.}
  \label{fig:63}
\end{figure}

For mesons propagation in $n_t$ and $(N_T - n_t)$ is identical up to a possible
relative minus sign and, when only the ground state is considered, the
correlator shows (depending on source and sink interpolators) a cosh- or
sinh-dependence on $n_t$:
\begin{equation}
  A_0\, e^{-n_t E_0} \pm A_0\, e^{-(N_T - n_t)E_0}
  = \begin{cases}
      2 A_0\, e^{-N_T E_0/2}\, \cosh\!\big( (N_T/2 - n_t) E_0 \big) \\
      2 A_0\, e^{-N_T E_0/2}\, \sinh\!\big( (N_T/2 - n_t) E_0 \big)
    \end{cases}
\end{equation}
In order to analyze in which range of $n_t$ the contribution of the sub-leading
exponentials in (6.54) can be neglected, one defines an effective mass as
\begin{equation}
  m_{\mathrm{eff}}\!\left( n_t + \half \right) = \ln
    \frac{C(n_t)}{C(n_t + 1)} .
\end{equation}
Once the correlator $C(n_t)$ is dominated by the ground state energy, $m_{\rm
eff}$ becomes constant and forms an effective mass plateau at $m_{\rm eff} =
E_0$.

If one wants to respect periodicity in $n_t$ according to (6.55), one sets
\begin{equation}
  \frac{C(n_t)}{C(n_t + 1)}
  = \frac{\cosh\!\big( m_{\rm eff}(n_t - N_T/2) \big)}
         {\cosh\!\big( m_{\rm eff}(n_t + 1 - N_T/2) \big)} ,
\end{equation}
(and equivalently for the case of sinh) and solves for $m_{\rm eff}$ at each
$n_t$. This was done in the right-hand side plot in Fig.~6.3. Depending on the
mass, such modifications due to periodicity effects may not be necessary if one
stays away from the region $n_t \lesssim N_T/2$.

In the right-hand side plot of Fig.~6.3 the effective mass plateaus are well
pronounced for $4 < n_t < 28$. In this range the contribution of excited states
can be neglected and a two-parameter fit to the simple form
$A_0 \exp(-n_t E_0)$ can be performed. If the correlator is symmetric, as in
our example, then one fits to $A_0 \cosh((n_t - N_T/2) E_0)$.

A generalization of effective mass techniques is combining correlation function
values from different time slices subtracting the leading exponential decay.
This may be extended to a tool for extraction excited hadron masses [18].

We stress that the symmetry of correlators, i.e., equal propagation in $n_t$ and
$N_T - n_t$ up to a sign, is not a general feature. In particular, after
projection to definite parity (see (6.16)) a baryon propagates in $n_t$, while
its parity partner propagates in $N_T - n_t$. Since the two masses are
different, the propagator is not symmetric, and instead of cosh or sinh, the
simple exponential form has to be used for effective masses and the fit
function.

\subsection{Fitting the correlators}

Although plotting $m_{\rm eff}$ gives a first estimate of the ground state
mass, a correlated fit in a range of $n_t$ values should be performed. For the
correlated analysis one fits the data in a range $n_{\min} \leq n_t, n_t' \leq
n_{\max}$ by minimizing
\begin{equation}
  \chi^2 = \sum_{n_t, n_t' = n_{\min}}^{n_{\max}}
    \big( C(n_t) - f(n_t) \big)\, w(n_t, n_t')\,
    \big( C(n_t') - f(n_t') \big) ,
\end{equation}
with regard to the parameters $A_0$, $E_0$ of the hypothesis function
\begin{equation}
  f(n_t) = A_0\, e^{-E_0 n_t} \qquad \text{or} \qquad
  f(n_t) = A_0\, \cosh\!\big( (n_t - N_T/2) E_0 \big) \ \text{(or sinh)} .
\end{equation}
The weight $w$ ideally should be the inverse of the exact covariance matrix,
i.e., $w(n_t, n_t') = \mathrm{Cov}^{-1}(n_t, n_t')$. In actual Monte Carlo runs
one only knows an estimator, i.e., the measured covariance matrix of the data,
\begin{equation}
  \mathrm{Cov}_N(n_t, n_t')
  = \frac{1}{N - 1}\Big\langle
      \big( C(n_t) - \langle C(n_t)\rangle_N \big)\,
      \big( C(n_t') - \langle C(n_t')\rangle_N \big)
    \Big\rangle_N ,
\end{equation}
where $\langle\ldots\rangle_N$ denotes the statistical average over $N$ gauge
configurations. In (6.58) one thus uses
$w(n_t, n_t') = \mathrm{Cov}^{-1}_N(n_t, n_t')$.

However, this estimator is often too badly determined and due to statistical
fluctuations there may be accidental small eigenvalues destabilizing the fit. In
such a case often either some smooth approximation or just the diagonal part of
the covariance matrix (neglecting the correlation between different data
points) are used. The latter gives rise to $w(n_t, n_t') =
\delta_{n_t,n_t'}/\sigma(n_t)^2$ and the $\chi^2$-functional of (6.58) reduces
to the form familiar from uncorrelated fits. For a critical account of
correlated fits, see [19].

One method to determine the fit range $[n_{\min}, n_{\max}]$ in (6.58) is the
analysis of effective mass plateaus discussed above. Another criterion is to
require that the $\chi^2$ per degree of freedom of the fit does not change by
more than $\order{1}$ when the region is extended. Although a generally valid
way to determine the fit interval is desirable, often -- in particular for data
with only moderate statistics -- it is affected by individual prejudice. Any
such decision involves some arbitrariness and should be checked by visual
inspection of the fits.

A simple method to estimate the statistical error of the fit result for the
mass is to use the statistical bootstrap or jackknife methods. One repeats the
analysis for large subsets of the data and determines the final error of the
mass a posteriori, as discussed in Chap.~4.

\subsection{The calculation of excited states}

In particular at smaller time distances one observes the contribution of states
with energies higher than that of the ground state (see the right-hand side
plot of Fig.~6.3). Obviously it is desirable to get an estimate of these energy
levels corresponding to excited states. At the same time, disentangling ground
states from higher states provides cleaner signals for both.

A direct fit of the correlation function to a hypothesis function $f(n_t)$
which is a sum of $K$ exponentials,
\begin{equation}
  f(n_t) = \sum_{k=1}^{K} c_k\, e^{-n_t E_k} ,
\end{equation}
would formally work, if the number of parameters were smaller than the number
of data points and if the data $C(n_t)$ were exact. In reality the data result
from Monte Carlo simulations and do have errors. It can be easily demonstrated
that already a sum of three exponentials may lead to unreliable results in such
fits. Several alternative approaches have been discussed in the literature.

\textbf{Bayesian analysis (conditionally biased fits):} In order to stabilize
the fit to a hypothesis function of the type (6.61), one does not use the
$\chi^2$ functional (6.58), but instead one minimizes
\begin{equation}
  F = \chi^2 + \lambda\, \phi ,
\end{equation}
where $\phi$ is some stabilizing function of the fit parameters and $\lambda$ a
positive real multiplier. This technique is well known for so-called ill-posed
problems like analytic continuation of functions given only by data points
[20].

In [21] it has been proposed to use as the stabilizing function
\begin{equation}
  \phi = \sum_{k=1}^{K}\left[ a_k\, (E_k - \bar{E}_k)^2
    + b_k\, (c_k - \bar{c}_k)^2 \right],
\end{equation}
where $\bar{E}_k$ and $\bar{c}_k$ are prejudices for the expected energy values
and coefficients. Adjusting the parameters $a_k$ and $b_k$ determines the
relative weight for the corresponding bias. The fit is thus stabilized, even if
the number of fit parameters exceeds the number of data, leading to unique
results depending on the multiplier $\lambda$, which defines the ``amount of
prejudice''. One tries to find regions for the resulting fit parameters that
depend only weakly on $\lambda$.

One technique based on such ideas [22] has been to start with simple
one-exponential fits at large-$n_t$ values, then continue with two-exponential
fits including the result of the earlier fit as a bias along the lines
discussed. Experience shows that one usually needs good statistics (small
statistical errors of the data) in order to successfully employ that method.

\textbf{Maximum entropy method:} For this approach one writes the correlation
function as the Laplace transform of a spectral density,
\begin{equation}
  C(n_t) = \int_0^\infty \dd E\, \rho(E)\, e^{-n_t E} .
\end{equation}
In the continuum the resulting spectral density $\rho(E)$ should exhibit peaks
near the energy values dominating the correlation function. When computing
$\rho(E)$ from lattice data, one discretizes the energy in small steps,
$E_n = n\,\Delta E$, and uses the values $\rho(E_n)$ of the density at the
discrete energies $E_n$ as fit parameters. The technical challenge is to
recover the spectral density from the (few) values of $C(n_t)$ by minimizing a
suitable distance functional (6.62). Various approaches to define the
stabilizing functional $\phi$ have been used.

In the maximum entropy method the functional $\phi$ in (6.62) is chosen based
on probabilistic arguments. We normalize the spectral density according to
\begin{equation}
  C(0) = \int_0^\infty \dd E\, \rho(E) \equiv 1 .
\end{equation}
Then $\rho(E)$ should be distributed such that each ``quantum'' has the same a
priori chance to contribute to the energy value $E$. This leads to the
stabilizer
\begin{equation}
  \phi[\rho] = \int_0^\infty \dd E\, \rho(E)\, \ln\!\big( \rho(E) \big) .
\end{equation}
Minimizing the joint probability leads to a smooth solution. A drawback is that
in this method the spectral density has to be positive. Also, the functional is
non-linear and dedicated numerical methods have to be used to solve the
minimization problem.

This approach is heavily used in statistical physics. It also has been applied
to nucleon propagators [23, 24]. From experience with statistical physics it is
expected, however, that one needs data for many values of $n_t$ and with very
small errors for a reliable result. Sample calculations show significant
ambiguities when one tries to extract more energy levels from moderately
precise correlation data.

\textbf{Variational analysis:} The situation concerning excited states can be
improved by considering not only a single correlator, but by computing a matrix
of cross correlators
\begin{equation}
  C_{ij}(n_t) \equiv \big\langle \widetilde{O}_i(0, n_t)\, O_j(0, 0) \big\rangle
  = \sum_k \bra{0}\hat{O}_i\ket{k}\,\bra{k}\hat{O}_j^\dagger\ket{0}\,
    e^{-n_t E_k} ,
\end{equation}
for a set of $N$ basis interpolators $O_i$, $i = 1, \ldots, N$, all with the
quantum numbers of the state one is interested in. Different interpolators can,
e.g., be constructed by using different Dirac structures as discussed in
Sect.~6.1, or with quark propagators acting on different quark sources.

Utilizing the special form of the spectral decomposition on the right-hand side
of (6.67) it has been shown [25--27] that diagonalizing the cross-correlation
matrix $C(n_t)$ allows to disentangle the physical states to some amount.
Indeed the eigenvalues $\lambda^{(k)}(n_t)$ of the correlation matrix can be
shown to behave as
\begin{equation}
  \lambda^{(k)}(n_t) \propto e^{-n_t E_k}\left( 1 +
    \order{e^{-n_t \Delta E_k}} \right),
\end{equation}
where $\Delta E_k$ is the distance of $E_k$ to nearby energy levels. Even
better results can be obtained when using the generalized eigenvalue equation
\begin{equation}
  C(n_t)\, v = \lambda(n_t)\, C(n_0)\, v ,
\end{equation}
where the eigenvalues behave again as in (6.68), but the amplitude of the
correction term is typically smaller. The generalized eigenvalue problem (6.69)
may be rewritten as
\begin{equation}
  Q(n_0)^{-1}\, C(n_t)\, Q^\dagger(n_0)^{-1}\, u = \lambda(n_t)\, u
  \qquad \text{with} \qquad Q(n_0)\, Q^\dagger(n_0) = C(n_0) .
\end{equation}
Numerically this form is expected to be more stable than diagonalizing
$C(n_0)^{-1}C(n_t)$, which is another way of rewriting the generalized
eigenvalue problem as a regular eigenvalue problem.

In the generalized eigenvalue problem the normalization at some time slice
$n_0 < n_t$ ought to improve the signal by suppressing the contributions from
higher excited states. Indeed that effect can be observed in sample calculations
using model data: Both approaches, the straightforward diagonalization, as well
as the generalized eigenvalue problem, give the correct answers asymptotically,
but the generalized eigenvalue approach does so already at smaller distances
$n_t$.

The method improves when the number of basis interpolators is increased. These
should be independent and have good overlap with the eigenstates of the
problem. On the other hand, in realistic calculations, including more
interpolators enhances the statistical noise and thus affects the
diagonalization. It always pays off to choose lattice operators which are close
to the expected physical content of the eigenstates.

The eigenvectors of the problem give the overlap of the eigenstates, supposedly
the states close to the physical states, with the basis interpolators used in
the correlation matrix (6.67). Thus, one may derive information on the wave
function of the physical modes. In [28--30] lattice operators have been
obtained by combining Jacobi smeared quark sources of different widths thus
allowing for nodes in the spatial wave function. A different approach is
followed in [31, 32] where extended lattice interpolators are constructed using
irreducible representations of the symmetry group of a cubic lattice.

\section{Finalizing the results for the hadron masses}

Now we are at a stage where the hadron masses have been extracted from the
corresponding correlators and are ready for analysis. We discuss the necessary
steps for finalizing the spectroscopy calculation and illustrate them using
data from quenched calculations.

\begin{figure}[htbp]
  \centering
  \includegraphics[width=0.9\textwidth]{fig64.png}
  \caption*{Fig. 6.4. \emph{Raw data for hadron masses. The left-hand side plot
    shows pion, $\rho$, and nucleon masses in lattice units as a function of the
    bare quark mass. The right-hand side plot is a so-called APE plot where
    dimensionless ratios are compared to the experimental number. The data are
    taken from the quenched calculation [33]}.}
  \label{fig:64}
\end{figure}

\subsection{Discussion of some raw data}

Following the steps of the last sections we have computed hadron masses in
lattice units, i.e., we have obtained numbers for the dimensionless products
$aM$. Usually, the hadron masses are evaluated for a range of light quark
masses $m_q \equiv m_u = m_d$. The left-hand side plot in Fig.~6.4 shows such
raw data for pion, $\rho$, and nucleon masses from a quenched simulation [33].

Converting masses in lattice units to numbers, that can be compared to
experimental data, can be done by considering ratios of masses where the
lattice constant $a$ cancels. The right-hand side plot in Fig.~6.4 shows
$M_N/M_\rho$, the ratio of nucleon to $\rho$ mass, as a function of
$(M_\pi/M_\rho)^2$, where $M_\pi$ is the pion mass. Such a plot is called APE
plot [34]. An alternative is the so-called Edinburgh plot [35], where one uses
$M_\pi/M_\rho$ on the horizontal axis.

In both, APE and Edinburgh plots, the data for $M_N/M_\rho$ can be compared to
the physical point, $M_N/M_\rho \approx 1.209$ at $M_\pi/M_\rho \approx 0.180$.
Although we do not perform an extrapolation, it is obvious that the data shown
in the APE plot of Fig.~6.4 clearly extrapolate to the physical point already
in the quenched approximation.

\subsection{Setting the scale and the quark mass parameters}

The results of a lattice calculation are always dimensionless numbers like,
e.g., mass ratios. In order to compare with the experiment, one has to relate a
mass or length unit to an experimental value. Consequently, the data have to be
converted to physical units such as MeV. In other words, one has to determine
the physical value of the lattice constant $a$ in order to obtain a mass $M$
from the calculated number for $aM$. There are two main techniques for this
step called scale setting.

\textbf{Setting the scale with the Sommer parameter:} For this approach one
uses the Sommer parameter to determine the lattice constant $a$ from the static
potential as discussed in Sect.~3.5. The result has the form $a = x$~fm, where
$x$ is a dimensionless number.

From our hadron correlators we have computed the dimensionless number
$X = aM$ for some hadron mass $M$. Consequently the physical mass $M = X/a$ is
given in inverse fm. This mass unit is a consequence of using natural units
$\hbar = 1$, $c = 1$. For converting this result to MeV one uses
\begin{equation}
  1 = \hbar c = 197.327\ \mathrm{MeV\,fm}
  \qquad\Longrightarrow\qquad
  1\ \mathrm{fm}^{-1} = 197.327\ \mathrm{MeV} .
\end{equation}
For the hadron mass in MeV we then obtain
\begin{equation}
  M = \frac{X}{a} = \frac{X}{x}\ \mathrm{fm}^{-1} = \frac{X}{x}\,
    197.327\ \mathrm{MeV} .
\end{equation}

\textbf{Setting the scale with a hadron mass:} In this approach one uses one of
the hadron masses $M_0$ for setting the scale. Popular candidates are the
$\rho$ meson or the nucleon. For that hadron, e.g., the $\rho$-meson, the mass
in lattice units, $aM_0(m_q)$, has to be known at several values of the bare
quark mass $m_q$. One then extrapolates the data to the value $m_q^*$ of the
mass parameter, where the dimensionless pion to $\rho$ mass ratio assumes its
physical value, $M_\pi(m_q^*)/M_\rho(m_q^*) = M_\pi^{\rm exp}/M_\rho^{\rm exp} =
0.180$. Alternatively one may also extrapolate $aM_0(m_q)$ to the chiral limit,
i.e., to $m_q = 0$. This is only an approximation for the true physical point,
but on the other hand this choice avoids the statistical error due to the
measured pion to $\rho$ mass ratio. For both extrapolations the result is a
number $X = aM_0(m_q^*)$ (or $X = aM_0(0)$), and we identify
$M_0(m_q^*) \equiv M_0^{\rm exp}$ (or $M_0(0) \equiv M_0^{\rm exp}$), where
$M_0^{\rm exp}$ is the experimental mass of our hadron in MeV. One obtains the
lattice constant in inverse MeV: $a = X/M_0^{\rm exp}$. This scale $a$ can be
directly used to obtain other masses $M$ in MeV from the values $aM$ in lattice
units.

For simulations with three quark flavors, in addition to the scale $a$, also
the light quark mass $m_q = m_u = m_d$ and the strange quark mass $m_s$ have to
be determined. For the light quark mass one either identifies the physical
point by setting $m_q = m_q^*$, with $m_q^*$ determined as discussed above, or
alternatively one assumes that the light $u$ and $d$ quarks are massless. For
setting the strange quark mass, a similar strategy is applied: One uses a
hadron containing the strange quark, typically the $K^\pm$ pseudoscalar meson,
the vector meson $\phi$, or baryons like $\Sigma$ or $\Omega$. The $K^\pm$ has
quark content $u$, $s$ and thus its mass is a function of the $u$ and $s$ quark
masses, $M_{K^\pm}(m_q, m_s)$ with $m_q = m_u$. In a first step the data are
extrapolated either to the chiral limit, $m_q = 0$, or to $m_q = m_q^*$. This
gives the $K^\pm$ mass $M_{K^\pm}(m_s)$ at different values of the strange quark
mass (the scale is set with one of the methods discussed above). This is done
for several values of $m_s$ and the data for $M_{K^\pm}(m_s)$ are interpolated.
The value $m_s^*$, where the interpolated data coincide with the experimental
value, $M_{K^\pm}(m_s^*) = M_{K^\pm}^{\rm exp} = 494$\,MeV, gives the bare
strange quark mass parameter $m_s = m_s^*$. When using the $\phi$, the situation
is even simpler, since $\phi$ is predominantly a strange--antistrange state and
no chiral extrapolation is needed.

\subsection{Various extrapolations}

Analyzing lattice data one faces the notorious, three-fold group of problems:
\begin{itemize}
  \item Finite size effects and infinite volume limit $V \to \infty$.
  \item Scaling in the continuum limit $a \to 0$.
  \item Chiral extrapolation $m \to 0$.
\end{itemize}

\textbf{Finite size:} All numerical lattice results are obtained in finite
volumes. At fixed other parameters one therefore has to study the dependence on
the lattice size and, if possible, extrapolate to infinite spatial volume. For
the extrapolation one may either use a naive parameterization, e.g., in powers
of $1/L$ for spatial extent $L$, or, if possible, a model based on theory.

Actually, like in the analysis of the spin models of statistical physics,
finite size effects may even be useful to learn about infinite volume results.
An example is chiral perturbation theory which provides a systematic expansion
with the infinite volume quantities appearing as parameters [36, 37]. We apply
such formulas in the next chapter when we discuss a lattice calculation of the
chiral condensate on finite lattices.

For most spectroscopy studies one uses $N^3 \times N_T$ lattices which are much
larger in the time direction than in the spatial directions and $L = aN$.
Indeed, neglecting the finite time extent (and thus finite temperature effects)
the leading finite size effects are due to the spatial volume. Effects from
interaction around the spatial torus lead to exponential corrections to the
mass of $\order{\exp(-\alpha L)}$ (cf.\ [38, 39] and the discussion in
Chap.~11). The leading contribution comes from the smallest mass hadron, the
pion, and is of $\order{\exp(-L M_\pi)}$. As a rule of thumb, for
$L M_\pi > 4$ finite size effects from this mechanism can be ignored.

For small volumes numerical studies of the volume dependence of hadron masses
seem to indicate a $1/L^n$ behavior, with $n \approx 2{-}3$. This would be
compatible with the intuitive picture that the wave function is squeezed at too
small spatial volume and that the observed shift in energy is in leading order
due to this squeezing effect (cf.\ [40, 41] and references therein). Most
analyses agree on the observation that ground state hadron masses obtained for
linear lattice sizes larger than 3~fm show no noticeable volume dependence. For
excited states with wave functions of larger extent, the squeezing effect may
be more serious.

\textbf{Continuum limit and scaling:} A measured mass $M$ (physical units) has
$a$-dependent corrections,
\begin{equation}
  M(a) = M_{\rm phys}\left( 1 + \order{a^\alpha} \right),
\end{equation}
which may be different for different lattice actions. Whereas with the Wilson
fermion action a quadratic contribution to the mass was identified [42], one
finds that improving the fermion action reduces this dependence [33].

Eventually we want to obtain results for continuous space--time. This
corresponds to the limit $a \to 0$, or, equivalently, $\beta \to \infty$.
However, decreasing the lattice spacing for a fixed number of lattice points
shrinks the physical volume. Thus, as we decrease $a$ (increase $\beta$), we
have to increase the number of lattice points, $N^3 \times N_T$, which in turn
drives up the numerical cost of the simulation. Typical simulations try to
afford three or four values of $a$ and perform an extrapolation to $a = 0$.

\textbf{Chiral extrapolation:} Often numerical results are obtained for
unphysically large quark masses. Eventually we want to approach the small quark
masses corresponding to the physical pion. From a conceptual point of view one
even might be interested in the massless case, the chiral limit $m \to 0$.
Massless 2-flavor QCD (i.e., assuming vanishing $m_u$ and $m_d$) is an
interesting and nontrivial, one-parameter theory, where hadron masses are
expected to be close to their experimental values. One could even call this a
zero-parameter theory, since the gauge coupling is only used to fix the scale.
Unfortunately this limit is hard to achieve for several reasons. Toward smaller
quark masses a Pandora's box of problems opens up.

For fermion actions that do not preserve chiral symmetry (or do not introduce a
regulator like the twisted mass action, cf.\ Chaps.~7 and 10), there are zero
modes occurring even for nonvanishing quark mass, due to the exceptional
configurations discussed in Sect.~6.2.5. This limits the smallest pion masses
that can be reached.

Another class of problems at small $m_q$ turns up when using the quenched
approximation. Omitting the fermion determinant has several effects:
\begin{itemize}
  \item The absence of the fermion determinant alters the physics by omitting
        sea quark loops (compare Fig.~6.2). Thus, hadrons cannot decay in the
        quenched approximation.
  \item The contribution of the anomaly (see next chapter) is removed and the
        $\eta'$-meson becomes light -- an additional Goldstone mode appears.
  \item States with negative coefficients in their correlation function appear,
        so-called ghosts. The quenched $\eta'$ combined with a hadron provides
        an example for such a state [43--45], cf.\ the hairpin diagrams in
        Fig.~6.2.
  \item Small eigenvalues of the Dirac operator are not suppressed by the
        determinant (which is the product of all eigenvalues). Consequently,
        fluctuations leading to exceptional configurations may occur more
        often.
\end{itemize}

For extrapolations in the quark mass, chiral perturbation theory provides the
suitable conceptional background (for introductory accounts, see [46--49]).
Chiral perturbation theory is an expansion in the pion mass around the chiral
limit taking into account the symmetries of QCD. The resulting extrapolation
formulas are typically expressed in terms of the pion mass, and as an example
we show the result for the extrapolation of the nucleon [50, 51]
\begin{equation}
  M_N = c_0 + c_2 M_\pi^2 + c_3 M_\pi^3 + c_4 M_\pi^4 \ln(M_\pi)
    + \order{M_\pi^4} ,
\end{equation}
where the expansion coefficients $c_i$ are related to low energy constants such
as the pion decay constant in the chiral limit.

Chiral perturbation theory can be modified to take into account quenching
[52--54]. Formulas such as (6.74) obtain corrections which may become singular
in the chiral limit. This is due to the additional Goldstone boson already
addressed.

\subsection{Some quenched results}

High precision determinations of hadron masses involve measurements on several
hundred gauge field configurations on lattices with size up to $64^3 \times
112$ for quenched Wilson fermions (see e.g.\ [42, 55]; simulations with other
formulations will be addressed in later chapters). The lattice spacing in these
simulations was as small as 0.05~fm at a linear physical lattice size of
$\order{3.2~\mathrm{fm}}$.

\begin{figure}[htbp]
  \centering
  \includegraphics[width=0.8\textwidth]{fig65.png}
  \caption*{Fig. 6.5. \emph{Quenched light hadron spectrum compared with
    experiments; statistical and systematic errors are indicated. The plot is
    from a high statistics simulation of the CP--PACS collaboration using Wilson
    valence quarks [42]. (Reprinted figure with permission from S.~Aoki et al.,
    Phys.\ Rev.\ D 67, 034503 (2003). Copyright (2003) by the American Physical
    Society)}.}
  \label{fig:65}
\end{figure}

To illustrate what has been achieved in quenched lattice spectroscopy, in
Fig.~6.5 we show the results of a high statistics study from the CP--PACS
collaboration [42]. The mass values have been obtained after extrapolation to
infinite volume, vanishing lattice spacing and the chiral limit. Results from
setting the strange quark mass with either $K$ or $\phi$ are compared. Although
the calculation was done in the quenched approximation, a good agreement with
experimental data is found for a large range of hadrons. In Chap.~8 we will
compare Fig.~6.5 to a more recent spectroscopy result with fully dynamical
quarks.

\section*{References}

\begin{enumerate}[label={[\arabic*]}]
  \item C.~Amsler et al.\ [Particle Data Group]: Phys.\ Lett.\ B 667, 1 (2008) 125
  \item D.B.~Leinweber et al.: in Lect.\ Notes Phys.\ 663, 71. Springer, Berlin,
        Heidelberg (2005) 129, 130
  \item S.~G\"usken et al.: Phys.\ Lett.\ B 227, 266 (1989) 138
  \item C.~Best et al.: Phys.\ Rev.\ D 56, 2743 (1997) 138
  \item R.~Barrett et al.: \emph{Templates for the Solution of Linear Systems:
        Building Blocks for Iterative Methods} (SIAM, Philadelphia 1994) 139,
        140, 141
  \item W.~H.~Press, S.~A.~Teukolsky, W.~T.~Vetterling, and B.~P.~Flannery:
        \emph{Numerical Recipes in C}, 2nd ed.\ (Cambridge University Press,
        Cambridge, New York 1999) 139
  \item H.~A.~van der Vorst: SIAM J.\ Sci.\ Statist.\ Comput.\ 13, 631 (1992)
        140
  \item M.~H.~Gutknecht: SIAM J.\ Sci.\ Statist.\ Comput.\ 14, 1020 (1993) 140,
        141
  \item I.~Montvay and G.~M\"unster: \emph{Quantum Fields on a Lattice}
        (Cambridge University Press, Cambridge, New York 1994) 141
  \item A.~Frommer et al.: Int.\ J.\ Mod.\ Phys.\ C 6, 627 (1995) 141
  \item B.~Jegerlehner: Nucl.\ Phys.\ B (Proc.\ Suppl.) 63, 958 (1998) 141
  \item M.~Albanese et al.: Phys.\ Lett.\ B 192, 163 (1987) 142
  \item A.~Hasenfratz and F.~Knechtli: Phys.\ Rev.\ D 64, 034504 (2001) 143
  \item C.~Morningstar and M.~Peardon: Phys.\ Rev.\ D 69, 054501 (2004) 143
  \item W.~Kamleh, D.~Adams, D.~B.~Leinweber, and A.~G.~Williams: Phys.\ Rev.\
        D 66, 014501 (2002) 143
  \item A.~Hasenfratz, R.~Hoffmann, and S.~Schaefer: JHEP 0705, 029 (2007) 143
  \item S.~D\"urr: Comput.\ Phys.\ Commun.\ 180, 1338 (2009) 143
  \item D.~Guadagnoli, M.~Papinutto, and S.~Simula: Phys.\ Lett.\ B 604, 74
        (2004) 146
  \item C.~Michael: Phys.\ Rev.\ D 49, 2616 (1994) 146
  \item S.~Ciulli, C.~Pomponiu, and I.~Sabba-Stefanescu: Phys.\ Rep.\ 17, 133
        (1975) 147
  \item G.~P.~Lepage et al.: Nucl.\ Phys.\ (Proc.\ Suppl.) 106, 12 (2002) 147
  \item Y.~Chen et al.\ (unpublished) arXiv:hep-lat/0405001 (2004) 148
  \item M.~Asakawa, T.~Hatsuda, and Y.~Nakahara: Prog.\ Part.\ Nucl.\ Phys.\ 46,
        459 (2001) 148
  \item K.~Sasaki, S.~Sasaki, T.~Hatsuda, and M.~Asakawa: Nucl.\ Phys.\ B
        (Proc.\ Suppl.) 129, 212 (2004) 148
  \item C.~Michael: Nucl.\ Phys.\ B 259, 58 (1985) 149
  \item M.~L\"uscher and U.~Wolff: Nucl.\ Phys.\ B 339, 222 (1990) 149
  \item B.~Blossier et al.: JHEP 0904, 094 (2009) 149
  \item T.~Burch et al.: Phys.\ Rev.\ D 70, 054502 (2004) 149
  \item T.~Burch et al.: Phys.\ Rev.\ D 73, 094505 (2006) 149
  \item T.~Burch et al.: Phys.\ Rev.\ D 74, 014504 (2006) 149
  \item S.~Basak et al.: Phys.\ Rev.\ D 72, 094506 (2005) 149
  \item S.~Basak et al.: Phys.\ Rev.\ D 72, 074501 (2005) 149
  \item C.~Gattringer et al.\ [BGR (Bern-Graz-Regensburg) collaboration]: Nucl.\
        Phys.\ B 677, 3 (2004) 150, 153
  \item P.~Bacilieri et al.: Phys.\ Lett.\ B 214, 115 (1988) 150
  \item K.~C.~Bowlers et al.: Phys.\ Lett.\ B 162, 354 (1985) 150
  \item J.~Gasser and H.~Leutwyler: Ann.\ Phys.\ 158, 142 (1984) 152
  \item P.~Hasenfratz and H.~Leutwyler: Nucl.\ Phys.\ B 343, 241 (1990) 152
  \item M.~L\"uscher: in \emph{Progress in Gauge Field Theory}, Carg\`ese, 1983,
        edited by G.~'t~Hooft et al.\ (Plenum, New York 1984) 152
  \item M.~L\"uscher: Commun.\ Math.\ Phys.\ 104, 177 (1986) 152
  \item M.~Fukugita et al.: Phys.\ Lett.\ B 294, 380 (1992) 152
  \item K.~Sasaki and S.~Sasaki: Phys.\ Rev.\ D 72, 034502 (2005) 152
  \item S.~Aoki et al.\ [CP-PACS collaboration]: Phys.\ Rev.\ D 67, 034503
        (2003) 153, 154, 155
  \item S.~R.~Sharpe: Phys.\ Rev.\ D 56, 7052 (1997) 153
  \item S.~R.~Sharpe: Phys.\ Rev.\ D 62, 099901 (2000) 153
  \item J.~Heitger, R.~Sommer, and H.~Wittig: Nucl.\ Phys.\ B 588, 377 (2000)
        153
  \item H.~Leutwyler: in \emph{Festschrift in Honor of B.L.~Ioffe: At the
        Frontier of Particle Physics}, edited by M.~Shifman, Vol.~1, p.~271
        (World Scientific, Singapore 2001) 154
  \item U.~Mei\ss ner: in \emph{Festschrift in Honor of B.L.~Ioffe: At the
        Frontier of Particle Physics}, edited by M.~Shifman, Vol.~1, p.~417
        (World Scientific, Singapore 2001) 154
  \item U.~Mei\ss ner: PoS LAT2005, 009 (2006) 154
  \item V.~Bernard and U.~Mei\ss ner: Ann.\ Rev.\ Nucl.\ Part.\ Sci.\ 57, 33
        (2007) 154
  \item T.~Becher and H.~Leutwyler: Eur.\ Phys.\ J.\ C 9, 643 (1999) 154
  \item M.~Procura, T.~R.~Hemmert, and W.~Weise: Phys.\ Rev.\ D 69, 034505
        (2004) 154
  \item S.~R.~Sharpe: Phys.\ Rev.\ D 46, 3146 (1992) 154
  \item C.~W.~Bernard and M.~F.~L.~Golterman: Phys.\ Rev.\ D 46, 853 (1992) 154
  \item C.~W.~Bernard and M.~F.~L.~Golterman: Phys.\ Rev.\ D 49, 486 (1994) 154
  \item K.~C.~Bowlers et al.\ [UKQCD collaboration]: Phys.\ Rev.\ D 62, 054506
        (2000) 154
\end{enumerate}
